\documentclass[11pt]{article}

\usepackage[margin=1in]{geometry}
\usepackage{amsmath,amssymb,amsthm,mathtools,mathrsfs}
\usepackage[hidelinks]{hyperref}

\newtheorem{theorem}{Theorem}
\newtheorem{lemma}[theorem]{Lemma}
\newtheorem{proposition}[theorem]{Proposition}
\newtheorem{corollary}[theorem]{Corollary}
\newtheorem{claim}[theorem]{Claim}
\theoremstyle{definition}
\newtheorem{definition}[theorem]{Definition}
\theoremstyle{remark}
\newtheorem{remark}[theorem]{Remark}

\newcommand{\I}{[0,1]}
\newcommand{\Fix}{\operatorname{Fix}}
\newcommand{\Crit}{\operatorname{Turn}}
\newcommand{\id}{\operatorname{id}}
\newcommand{\cl}{\operatorname{cl}}

\title{Common Fixed Points for Commuting Definable\\
Piecewise-\(C^1\) Interval Maps with Finitely Many Pieces\\
\large Formal-proof-aligned revision}
\author{}
\date{}

\begin{document}
\maketitle

\begin{abstract}
Let $P,Q:[0,1]\to[0,1]$ be commuting continuous maps which are
finite-piece piecewise $C^1$ and definable in a common o-minimal expansion of
the ordered real field.  We prove that $P$ and $Q$ have a common fixed point.
The class includes finite-piece piecewise-linear maps, finite-piece
piecewise-real-analytic maps, and branches with definable infinite-order
flatness.  The proof reduces a hypothetical counterexample by finite
fixed-component quotients, analyzes a repelling one-sided germ through its
orbit circle, obtains exactness, conjugates the centralizer to constant-slope
piecewise-linear maps using the Parry--Shultz scaling measure, and finishes
with a common Markov partition, Durand's substitution form of Cobham's
theorem, and an affine-circle reflection argument.
\end{abstract}

\begin{remark}[Status and proof contract]
This revision follows the dependency order and hypotheses of the Lean
formalization in \texttt{CommutingInterval}.  In particular, every quotient
is accompanied by an explicit fibre formula and lifting statement; every
one-sided germ records its nondegenerate domain; every use of an expanding
scale records the strict inequality (1<s); and the finite-arrival argument
uses an interior fixed barrier.  The research-level inputs used below are
stated in their exact consumed forms: one-variable o-minimal monotonicity and
bounded variation, local (C^{1+BV}) Kopell, the lift form of Denjoy,
Fugate--McLean, Parry--Shultz, and Durand--Cobham.  All remaining reductions
are part of the proof.
\end{remark}

\section*{Introduction}
Fix an o-minimal expansion $\mathcal R$ of the ordered real field; parameters
are allowed in the word ``definable.''  Thus every $\mathcal R$-definable
subset of the real line is a finite union of points and intervals.  A
continuous map $T\colon\I\to\I$ will be called
\emph{$\mathcal R$-definable finite-piece piecewise $C^1$} if it is
$\mathcal R$-definable and there is a finite partition
\[
 0=x_0<x_1<\cdots <x_n=1
\]
such that, on every closed interval $[x_{i-1},x_i]$, the map $T$ agrees
with an $\mathcal R$-definable $C^1$ function defined on an open
neighborhood of that interval.  Constant branches and derivative jumps at
partition points are allowed.  Both maps in the theorem are required to be
definable in the same o-minimal structure.

\begin{theorem}\label{thm:tame-C1}
Let $P,Q\colon\I\to\I$ be continuous $\mathcal R$-definable finite-piece
piecewise $C^1$ maps.  If $PQ=QP$, then $P$ and $Q$ have a common fixed
point.
\end{theorem}

\begin{corollary}\label{cor:analytic-PL}
The conclusion holds in particular for commuting continuous finite-piece
piecewise-linear maps and for commuting continuous finite-piece
piecewise-real-analytic maps of $[0,1]$, where the latter phrase means that
each analytic branch extends real-analytically to an open neighborhood of
its closed partition interval.
\end{corollary}

\begin{proof}
Finite-piece piecewise-linear maps are semialgebraic.  If a real-analytic
branch is defined on a neighborhood of a compact interval, an affine change
of variables realizes it as a restricted analytic function; finitely many
such branches are therefore definable in $\mathbb R_{\mathrm{an}}$.  Apply
Theorem~\ref{thm:tame-C1}.
\end{proof}

The unrestricted common-fixed-point assertion for commuting continuous
interval maps is false: Boyce and Huneke independently constructed
counterexamples \cite{Boyce,Huneke}.  Classical positive results under
extra hypotheses include Cohen's full-map theorem and the refinements of
Folkman and Joichi requiring only one member of the commuting pair to be
full \cite{Cohen,Folkman,Joichi}.  In a related finite-complexity direction,
Anu\v{s}i\'c and Mouron proved that strongly commuting piecewise-monotone
interval maps have a common fixed point \cite{AnusicMouron}; their
set-valued strong-commutation condition is strictly stronger than ordinary
commutation.  More precisely, if inverse images are regarded as set-valued
maps, ordinary commutation gives only
\(P\circ Q^{-1}\subseteq Q^{-1}\circ P\), whereas strong commutation requires
equality.  Thus even the piecewise-linear case of
Corollary~\ref{cor:analytic-PL} is not a formal consequence of that theorem:
here the extra equality is replaced by definable finite-piece
piecewise-\(C^1\) regularity.  Our conclusion is the existence of a common
fixed point, not a classification of commuting pairs.  In particular,
Corollary~\ref{cor:analytic-PL} covers the finite-piece piecewise-linear and
piecewise-real-analytic classes.
The real-analytic clause is branchwise: it allows turns and critical points
and does not assume that the branches assemble into a holomorphic selfmap
of a common complex domain.  It is therefore a statement about real
interval dynamics, rather than an application of a holomorphic fixed-point
theorem.

The definability assumption is a finite-geometric-complexity condition: it
rules out infinitely many connected components of zero sets, critical sets,
and level sets.  It is weaker than analyticity.  For example, the function
which is zero for $x\leq0$ and $e^{-1/x}$ for $x>0$ is definable in
$\mathbb R_{\mathrm{an},\exp}$.  No monotonicity, injectivity, or
nonflatness of either map is assumed.

Our convention is \(PQ=P\circ Q\).  Since all products below involve
commuting maps, their order causes no ambiguity.  The proof uses the
one-dimensional o-minimal monotonicity theorem, the bounded-variation form of
Kopell's lemma, the Poincar\'e--Denjoy circle theorems, the
Parry--Shultz scaled-measure theorem, Durand's substitution form of Cobham's
theorem, and the Fugate--McLean theorem on periodic homeomorphisms of
tree-like continua.  The reductions to the hypotheses of these results are
included below.  No theorem about commuting piecewise-linear interval maps is
used as a black box; the polygonal terminal argument is part of the proof.

Here is a roadmap.  Sections~1--3 pass to a minimal invariant core, make only
quotients respected by all maps in play, normalize a least counterexample,
and obtain local eventual surjectivity from a selected repelling side.
Section~4 replaces that one-sided dynamics by an orbit circle.  Kopell's
lemma eliminates rational rotation data, while Denjoy's theorem linearizes
the irrational case.  Sections~5--6 build the canonical ray, analyze its
tail core, and promote dense tails to exactness.  Section~7 uses the
Parry--Shultz measure to conjugate the relevant common iterate to constant
absolute slope.  Sections~8--9 construct a common finite Markov partition
and two substitutions for the same directed edge word; Durand's theorem
forces that word, hence the metric ray, to be eventually periodic.  Finally,
Section~10 globalizes the original phases to affine circle lifts and uses the
unique reflection together with a parity argument to produce a common fixed
point.

\par\noindent\textit{Proof of Theorem~\ref{thm:tame-C1}.}\quad
Suppose, toward a contradiction, that \(P,Q\) commute and have no common
fixed point.

\section{The tame piecewise-\texorpdfstring{\(C^1\)}{C1} category and safe quotients}

\begin{lemma}[Tame \(C^1\) finiteness]\label{lem:analytic-finiteness}
The class of \(\mathcal R\)-definable finite-piece piecewise \(C^1\) interval
maps has the
following properties.
\begin{enumerate}
\item It is closed under composition, restriction to a compact interval, and
affine changes of coordinates.
\item For every map \(T\) in the class, \(\Fix(T)\) is a finite union of
points and compact intervals.
\item If \(T\) has no flat interval, then it is finite-to-one and has finitely
many strict monotonicity laps.
\item On every displayed \(C^1\) branch, the derivative has bounded
variation.
\item Suppose finitely many pairwise disjoint compact intervals are collapsed
by the canonical finite piecewise-affine nondecreasing quotient, with positive
affine slope on every complementary gap, and a finite family of maps in the
class respects the quotient fibres.  The descended maps are again
\(\mathcal R\)-definable finite-piece piecewise \(C^1\).
\end{enumerate}
\end{lemma}

\begin{proof}
For composition, let \(U,T\) belong to the class.  The preimage under \(T\)
of each breakpoint of \(U\) is definable, hence is a finite union of points
and intervals.  Refine the source partition at the finitely many endpoints of
these components.  A resulting cell which is mapped constantly to a
breakpoint gives a constant composite branch.  Every other cell is mapped
into one \(C^1\) branch of \(U\), so \(U\circ T\) has a \(C^1\) extension
across that closed cell.  Definability is closed under composition.
Restriction and affine conjugacy are immediate.

The set \(\Fix(T)\), being the zero set of the definable function
\(T-\id\), is a finite union of points and intervals.  This proves (2).
The one-dimensional o-minimal monotonicity theorem
\cite{vanDenDries} says that a definable
one-variable function has a finite partition on whose interval cells it is
constant, strictly increasing, or strictly decreasing.  Apply it to \(T\).
If \(T\) has no flat interval, all its nondegenerate cells are strict, giving
finitely many strict laps and a uniform finite bound on its fibres.  This is
(3).

On a displayed branch let \(F\) be the definable \(C^1\) extension.  Its
derivative is definable: the graph of \(F'\) is described by the usual
first-order epsilon--delta formula.  Apply o-minimal monotonicity to the
continuous function \(F'\) on the compact branch.  On a finite partition it
is constant or monotone, so its total variation is the finite sum of the
absolute changes over those cells.  This proves (4).

For (5), choose the quotient map \(\pi\) to be constant on the collapsed
intervals and affine with positive slope on each complementary gap.  Put
\(H=\pi\circ T\), which belongs to the class by (1).  On the closure \(G\)
of each uncollapsed source gap, \(\pi|_G\) has an affine inverse \(s\) which
extends across the closed interval \(\pi(G)\).  Intersect the branch
intervals of \(H\) with \(G\).  On every resulting interval \(K\), the
descended map has the form
\[
  (\text{a }C^1\text{ extension of }H|_K)\circ s,
\]
and therefore has a definable \(C^1\) extension across \(\pi(K)\).  At a
collapsed source value, respect of quotient fibres says that \(H\) is
constant on the collapsed interval.  The extensions from the adjacent
noncollapsed gaps therefore take the same value there; adjoin the collapsed
point to those neighboring closed branch intervals (with the evident
one-sided convention at an ambient endpoint).  The descended graph
\(\{(\pi(x),\pi(T(x))):x\in\I\}\) is definable, and only finitely many pieces
occur.  This proves (5).
\end{proof}

Here and below, an interval \emph{invariant under \(T\)} is forward
invariant, \(T(J)\subseteq J\), unless equality is explicitly asserted.  Let
$\mathscr J$ be the family of nonempty compact intervals invariant under
both $P$ and $Q$, ordered by inclusion.  Every chain in $\mathscr J$ has a
nonempty compact intersection, and that intersection is again invariant.
Thus Zorn's lemma gives an inclusion-minimal member.  Restrict to it and
rescale affinely to $\I$.  It is nondegenerate, since a singleton would be a
common fixed point.  Moreover
\begin{equation}\label{eq:onto}
 P(\I)=\I=Q(\I).
\end{equation}
Indeed, $P(\I)$ and $Q(\I)$ are nonempty compact intervals invariant under
both maps, so minimality forces equality.  Henceforth there is no proper
nonempty compact interval invariant under both maps.

\begin{proposition}[Finite respected quotient package]
\label{prop:quotient-package}
Let $\pi:\I\to I'$ be the canonical nondecreasing quotient which collapses
finitely many pairwise disjoint compact intervals and is affine with positive
slope on every complementary gap.  Let $\mathscr F$ be a finite commuting
family of maps in the tame piecewise-$C^1$ class, and suppose every member of
$\mathscr F$ respects the fibres of $\pi$.  Then:
\begin{enumerate}
\item every $F\in\mathscr F$ descends to a unique continuous map
$\overline F:I'\to I'$ in the same tame class; the descended maps commute,
and surjectivity descends;
\item for every map $T$ respecting the fibres,
\begin{equation}\label{eq:fix-quotient-analytic}
 \Fix(\overline T)=\pi(\Fix(T));
\end{equation}
\item the inverse image of an interval invariant under all descended maps is
an interval invariant under all original maps;
\item if $T=P^rQ^s$ with $r,s>0$ and the collapsed intervals are
nondegenerate components of $\Fix(T)$, then $P$ and $Q$ respect the quotient,
and every common fixed point of the descended pair lifts to a common fixed
point of $P,Q$;
\item after fixed sets have become finite, a respected quotient followed by
restriction to a common invariant core cannot increase their cardinalities;
it strictly decreases the cardinality of $\Fix(T)$ whenever one quotient
fibre contains two distinct $T$-fixed points.
\item Suppose $P,Q\in\mathscr F$ have no common fixed point and admit no
proper nonempty compact interval which is forward invariant under both maps.
If $\pi$ is nonconstant, then the descended pair
$\overline P,\overline Q$ also has no common fixed point.
\end{enumerate}
\end{proposition}

\begin{proof}
The existence, continuity, and tame regularity of the descended maps are
Lemma~\ref{lem:analytic-finiteness}(5).  Uniqueness follows from the
surjectivity of $\pi$, and commutation follows after composing with $\pi$.
If $F$ is onto and $y=\pi(x)$, choose $z$ with $F(z)=x$; then
$\overline F(\pi(z))=y$, so $\overline F$ is onto.

One inclusion in \eqref{eq:fix-quotient-analytic} is immediate.  Conversely,
if $\overline T(\pi(x))=\pi(x)$, the fibre
$C=\pi^{-1}(\pi(x))$ is a compact interval and $T(C)\subseteq C$.
The interval fixed-point theorem supplies a point of $C\cap\Fix(T)$.
This proves (2).  If $J'\subseteq I'$ is invariant under every descended
map, then $\pi^{-1}(J')$ is a compact interval and the intertwining identity
$\pi F=\overline F\pi$ proves (3).

For (4), put $X=\Fix(T)$.  On $X$ the restrictions of the factors have the
explicit two-sided inverses
\[
 (P|_X)^{-1}=P^{r-1}Q^s|_X,\qquad
 (Q|_X)^{-1}=Q^{s-1}P^r|_X.
\]
Thus $P$ and $Q$ are homeomorphisms of $X$ and permute its connected
components.  They consequently respect the quotient which collapses the
nondegenerate components.  A common fixed quotient fibre is either a
singleton, in which case its point is fixed by both maps, or a component
$C\subseteq X$ with $P(C)=C=Q(C)$.  On $C$ the restrictions are commuting
interval homeomorphisms.  If both preserve orientation, both endpoints are
common fixed points.  If one reverses orientation, it has a unique fixed
point, and commutation forces the other map to fix that point.

Finally, \eqref{eq:fix-quotient-analytic} says that a quotient only identifies
fixed points, while restriction to a core only discards them.  If a fibre
contains two distinct fixed points, their images coalesce, giving the strict
inequality in (5).

For (6), suppose that $y\in I'$ is fixed by both descended maps and put
$C=\pi^{-1}(y)$.  The intertwining identities give
$P(C)\subseteq C$ and $Q(C)\subseteq C$.  If $C$ is a singleton, its point is
fixed by both original maps.  If $C$ is nondegenerate, then it is a proper
compact interval because $\pi$ is nonconstant, and it is forward invariant
under both maps.  Either alternative contradicts the hypotheses.
\end{proof}

All quotient operations below are canonical quotients of the preceding
form.  Proposition~\ref{prop:quotient-package} will be used without further
comment to retain the tame category, preserve a counterexample after a
nonconstant respected quotient of a minimal core, compare fixed-set
cardinalities, and lift a common fixed point from a fixed-component quotient.

\section{Normalization and the selected repelling side}

Put
\begin{equation}\label{eq:analytic-W}
 W=(PQ)^2.
\end{equation}

\begin{lemma}[Normalization and lexicographic descent]
\label{lem:analytic-normalization}
Starting from a tame piecewise-\(C^1\) counterexample on a minimal common
invariant core, one can obtain another such counterexample for which
\(X=\Fix(W)\) is finite and, on putting
\[
 \nu=\#X
\]
and
\begin{equation}\label{eq:analytic-words}
 L=2(2\nu)!,\qquad A=P^L,\qquad B=Q^L,\qquad
 R=AB=W^{L/2},\qquad S=R^2=W^L,
\end{equation}
the sets \(\Fix(R)\) and \(\Fix(S)\) are finite as well.  Normalized tame
counterexamples therefore exist, and among all of them one can choose a pair
minimizing
\begin{equation}\label{eq:analytic-mintriple}
 \bigl(\#\Fix(W),\#\Fix(R),\#\Fix(S)\bigr)
\end{equation}
lexicographically.

For such a minimizing pair, any later nonconstant quotient by finitely many respected
interval classes, followed by restriction to a minimal common invariant core
and the same fixed-component normalization, has nonincreasing fixed counts.
If its first count remains \(\nu\), its canonical exponent is still \(L\), so
the descended words \(R,S\) are literally those in
\eqref{eq:analytic-words}.
\end{lemma}

\begin{proof}
Lemma~\ref{lem:analytic-finiteness} says that \(\Fix(W)\) has finitely many
components.  If \(\Fix(W)=\I\), collapsing its sole nondegenerate component
gives the one-point quotient, whose unique point is fixed by both descended
maps; Proposition~\ref{prop:quotient-package}(4) would then lift a common
fixed point of \(P,Q\).  Hence this case is impossible, and the quotient
which collapses all nondegenerate components of \(\Fix(W)\) is nonconstant.
Collapse them and return to a minimal common invariant core.  The quotient
remains tame piecewise \(C^1\), a common fixed point would lift, and
\eqref{eq:fix-quotient-analytic} says that the new \(\Fix(W)\) is finite.
Every later respected quotient or core
restriction can only identify or discard points of this finite set: by
\eqref{eq:fix-quotient-analytic} it cannot increase its cardinality or
recreate a nondegenerate fixed component.

Define \(\nu,L,A,B,R,S\) as in the statement and perform the following
finite descent.  First consider \(\Fix(R)\).  It cannot equal \(\I\):
collapsing that single fixed component would give the one-point quotient,
and Proposition~\ref{prop:quotient-package}(4) would lift its common fixed
point.  Hence the fixed-component quotient for \(R\) is nonconstant.  Take
that quotient and return to a minimal core.  If the cardinality of
\(\Fix(W)\) has fallen, rename the descended pair \(P,Q\), recompute
\(\nu,L,A,B,R,S\), and begin this paragraph again.  Otherwise \(\nu\), and
therefore \(L\), is unchanged; \(R=P^LQ^L\) descends literally, and
\eqref{eq:fix-quotient-analytic} makes its fixed set finite.

Next apply the same operation to \(\Fix(S)\).  Again \(\Fix(S)=\I\) is
excluded by Proposition~\ref{prop:quotient-package}(4).  Collapse its
nondegenerate components and return to a minimal core.  If
\(\#\Fix(W)\) falls, rename the pair, recompute the six displayed objects,
and restart with \(R\).  If it does not fall, \(L\) is still unchanged,
\(\Fix(S)\) is now finite, and the already finite sets \(\Fix(W)\) and
\(\Fix(R)\) remain finite by
\eqref{eq:fix-quotient-analytic} and restriction.  Each restart strictly
lowers the positive integer \(\nu\), so only finitely many restarts occur.
The procedure therefore terminates with all three fixed sets finite, proving
that normalized tame counterexamples exist.

The triples in \eqref{eq:analytic-mintriple} form a nonempty subset of
\(\mathbb N^3\), so they have a lexicographic minimum.  Under a later
respected quotient and core restriction, the first fixed count cannot
increase.  A strict decrease already contradicts minimality, regardless of
the newly recomputed exponent.  If the first count is unchanged, the formula
\(L=2(2\nu)!\) gives the same exponent, so \(R,S\) descend literally and
their fixed counts also cannot increase.  Any further fixed-component
normalization only decreases those counts.  This proves the assertion.
\end{proof}

Fix henceforth a normalized counterexample minimizing
\eqref{eq:analytic-mintriple}.

The set \(X\) has at least two points.  Otherwise its unique point is
preserved by both factors and is common fixed.  On \(X\), the factors are
permutations, with
\begin{equation}\label{eq:analytic-X-inverses}
 (P|_X)^{-1}=Q(PQ)|_X,\qquad
 (Q|_X)^{-1}=P(PQ)|_X.
\end{equation}

Choose consecutive points \(p<c\) of \(X\).  The continuous function
\(W-\id\) has a strict constant sign on \((p,c)\).  Reflecting the ambient
coordinate if necessary, assume
\begin{equation}\label{eq:analytic-gap-sign}
 W(x)>x\qquad(p<x<c).
\end{equation}

Call a directed side based at a point of \(X\) \emph{\(W\)-repelling} if,
on some one-sided neighborhood, \(W\) preserves the side and moves every
point strictly away from its base.  The selected side \((p,+)\) is such a
side.  There are at most \(2\nu\) such states.

\begin{lemma}[Directed-state return]\label{lem:analytic-states}
The maps \(P,Q\) permute the finite set of \(W\)-repelling directed states.
Consequently \(A,B\) fix the selected state \((p,+)\), and their selected
germs are orientation-preserving \(C^1\) local diffeomorphisms.  Their
derivatives have bounded variation on a sufficiently small selected-side
interval.  They also
fix \(p\) and \(c\).
\end{lemma}

\begin{proof}
Put \(C_P=Q(PQ)\), so that \(C_PP=W\).  On a repelling state the one-sided
derivative of \(W\) is nonzero: in the coordinate measuring distance into
the side, \(W(t)>t\), and therefore
\[
 W'(0)=\lim_{t\downarrow0}\frac{W(t)}t\geq1.
\]
Before applying the chain rule, shrink the source side as follows.  If
\(\gamma(t)\), \(t\geq0\), is the directed source coordinate with
\(\gamma(0)=x\), the definable difference
\(P(\gamma(t))-P(x)\) is either identically zero or has a constant nonzero
sign on a sufficiently small punctured side.  The first alternative would make
\(W=C_PP\) constant there, so it is impossible.  The second selects one
directed side at \(P(x)\); shrink once more so that a single displayed
\(C^1\) branch of \(P\), and the corresponding directed branch of \(C_P\),
are used throughout.  The one-sided chain rule for \(C_PP=W\) now makes
both derivatives nonzero.  Thus the two selected germs are local
diffeomorphisms.  In directed distance coordinates the germ induced by
\(P\) is increasing, even if \(P\) reverses ambient orientation.  It maps
the state to the selected directed side at \(P(x)\in X\).  From \(PW=WP\),
this target germ is locally conjugate to the source germ, and hence is again
repelling.
Moreover \(C_P\) takes the target state back to the source because
\(C_PP=W\) and \(W\) preserves the source state.  Thus \(P\) is injective,
and hence a permutation, on the finite state set.  The same proof, using
\(C_Q=P(PQ)\), applies to \(Q\).

Every permutation of at most \(2\nu\) objects has order dividing
\((2\nu)!\).  Hence \(L\) fixes every directed state; a return to the same
directed side has positive orientation.  The same exponent fixes all points
of \(X\), since \(P,Q\) permute \(X\).  In particular \(A,B\) fix \(p,c\).
\end{proof}

\section{Finite arrival and local eventual surjectivity}

Put \(J=[p,c]\).  Since \(A,B\) fix its endpoints, continuity gives
\begin{equation}\label{eq:analytic-J-expands}
 J\subseteq A(J),\qquad J\subseteq B(J).
\end{equation}
Let
\[
 K_{i,j}=A^iB^j(J),\qquad
 K=\cl\bigcup_{n\geq0}K_{n,n}.
\]
The family is increasing in each index, and cofinality gives
\begin{equation}\label{eq:analytic-K-invariant}
 A(K)=K=B(K).
\end{equation}

We record the phase-intersection fact which controls every interval quotient
below.  If a compact interval \(C\) satisfies \(P^L(C)=C=Q^L(C)\), put
\[
 D_{i,j}=P^iQ^j(C),\qquad (i,j)\in(\mathbb Z/L\mathbb Z)^2.
\]
Translation of either label is an automorphism of the labelled intersection
graph.  One direction follows by applying \(P\) or \(Q\).  Conversely, if
\(y\in P(D_\alpha)\cap P(D_\beta)\), then
\[
 P^{L-1}(y)\in P^L(D_\alpha)\cap P^L(D_\beta)
 =D_\alpha\cap D_\beta,
\]
because \(P^L\) fixes every phase interval.  The argument for \(Q\) is
identical.  Hence the graph is vertex transitive, without any injectivity
assumption on \(P,Q\).

Apply this fact to \(C=K\), and set
\[
 F=\bigcup_{0\leq i,j<L}P^iQ^j(K).
\]
This is a finite union of compact intervals and \(P(F)=F=Q(F)\).  Every phase
interval contains the two distinct points \(P^iQ^j(p),P^iQ^j(c)\) of
\(X\), because the factors are injective on \(X\).

If \(F\) is disconnected, its connected components are finitely many
pairwise disjoint proper nondegenerate intervals permuted by \(P,Q\).  Indeed,
a continuous image of a component lies in one component; because
\(P(F)=F=Q(F)\), each induced selfmap of the finite component set is
surjective and hence bijective, and the unique source component maps onto its
target.  Each component contains at least two \(W\)-fixed points.  Collapse
these components.
Lemma~\ref{lem:analytic-finiteness} makes this a safe tame quotient,
Proposition~\ref{prop:quotient-package}(6) keeps the descended pair a
counterexample, and
\eqref{eq:fix-quotient-analytic} strictly lowers \(\#\Fix(W)\), contrary to
\eqref{eq:analytic-mintriple}.  Thus \(F\) is connected.  It is a common
invariant interval, so minimality forces \(F=\I\).

The phase-intersection graph is now connected and vertex transitive.  A
connected vertex-transitive interval graph is complete: an interval with
least right endpoint is simplicial, vertex transitivity makes every vertex
simplicial, and a shortest path rules out two nonadjacent vertices.  The
Helly property for intervals therefore gives
\[
 C_0=\bigcap_{0\leq i,j<L}P^iQ^j(K)\neq\varnothing.
\]
Label translation shows \(P(C_0)\subseteq C_0\) and
\(Q(C_0)\subseteq C_0\).  This intersection can be neither a singleton nor
a proper nondegenerate interval, so \(C_0=\I\).  In particular \(K=\I\).
Consequently
\begin{equation}\label{eq:analytic-dense-diagonal}
 R^n(J)\quad(n\geq0)
\end{equation}
is an increasing sequence of intervals with dense union in \(\I\).

Write \(R^n(J)=[\ell_n,r_n]\).  Then
\(\ell_n\downarrow0\) and \(r_n\uparrow1\).  We next prove that both ambient
endpoints occur at a finite stage.

Suppose, for definiteness, that \(0\) never occurs.  We claim that
\begin{equation}\label{eq:analytic-singleton-endpoint}
 S(0)=0,\qquad S^{-1}(0)=\{0\}.
\end{equation}
Every interior point belongs to some even diagonal interval
\(R^{2n}(J)\).  If such a point mapped to zero under \(S=R^2\), then zero
would belong to \(R^{2n+2}(J)\), contrary to the assumption.  If the other
ambient endpoint has already occurred, the same argument, starting from a
diagonal interval containing it, excludes it as a preimage of zero.  Since
\(S\) is onto, zero must then be its own unique preimage.

It remains to justify the case in which both ambient endpoints are
asymptotic.  Every \(x\in(0,1)\) lies in some diagonal interval contained in
\((0,1)\).  If \(R(x)\) were an ambient endpoint, that endpoint would occur
at the next stage.  Hence \(R((0,1))\subseteq(0,1)\).  Onto-ness now forces
the preimages of the two ambient endpoints to lie among the endpoints
themselves, so \(R\) permutes \(\{0,1\}\).  Thus \(S=R^2\) fixes both
endpoints and neither has any other preimage.  This proves
\eqref{eq:analytic-singleton-endpoint} in all cases.

On the selected \(C^1\) branch at zero, the derivative \(S'\) is definable.
Its zero set therefore has only finitely many components.  It cannot contain
an interval adjacent to zero, since that would make \(S\) constant there and
contradict \eqref{eq:analytic-singleton-endpoint}.  After shrinking
\(\varepsilon>0\), the derivative consequently has a constant nonzero sign
on \((0,\varepsilon)\).  Because \(S(0)=0\) and \(S(x)>0\) for \(x>0\), that
sign is positive.  Hence \(S\) is strictly increasing on
\([0,\varepsilon]\).
Compact separation gives
\[
 \eta=\min S([\varepsilon,1])>0.
\]
For all sufficiently large \(n\), both
\(\ell_{2n}<\varepsilon\) and \(\ell_{2n+2}<\eta\).  Choose
\(x_n\in R^{2n}(J)=[\ell_{2n},r_{2n}]\) with
\(S(x_n)=\ell_{2n+2}\).  The definition of \(\eta\) forces
\(x_n<\varepsilon\).  Since \(S\) is increasing on
\([0,\varepsilon]\) and \(\ell_{2n}\leq x_n\), we have
\(S(\ell_{2n})\leq\ell_{2n+2}\).  The reverse inequality follows from
\(S(\ell_{2n})\in S(R^{2n}(J))=R^{2n+2}(J)\).  Hence
\[
 \ell_{2n+2}=S(\ell_{2n}).
\]
The positive sequence on the left decreases to zero.  Since \(\Fix(S)\) is
finite, the strict sign of \(S(x)-x\) on a sufficiently small punctured
endpoint interval is therefore negative.  Thus every sufficiently small
positive point has iterates decreasing to zero.  More explicitly, after
shrinking \(\rho>0\), singleton-fibre positivity and the strict sign give
\[
 0<S(x)<x\qquad(0<x\leq\rho).
\]
Every orbit beginning in \((0,\rho]\) decreases to a fixed point in
\([0,\rho]\), necessarily zero.  If an arbitrary basin point first enters
\((0,\rho)\) after \(n\) iterates, continuity of \(S^n\) gives a neighborhood
which also enters \((0,\rho)\).  Hence the basin of zero is open.

Let \(U_0\) be the basin component containing zero and put
\(C_e=\cl(U_0)\).  The barrier needs a little care.  Phase-zero nonarrival
says (0\notin J=[p,c]); hence (0<p<c\leq1).  Thus (p), which is fixed by
(R) and therefore by (S=R^2), is an interior fixed point outside the
endpoint basin component.  Consequently (C_e) is a proper nondegenerate
endpoint interval.  (For a missing right endpoint the symmetric argument
uses (c<1) and the interior fixed barrier (c).)  If \(u\) is its free
endpoint, connectedness gives \(S(C_e)\subseteq C_e\).  The point \(u\) is
not in the basin.  Its component is the relative interval
\(U_0=[0,u)\), so if \(S(u)\in U_0\)---including \(S(u)=0\)---then \(u\)
would itself be in the basin.  Hence \(S(u)=u\), and
\begin{equation}\label{eq:analytic-basin-onto}
 S(C_e)=C_e.
\end{equation}

We also need the elementary singleton-fibre factorization.  If commuting
onto maps \(U,V\) satisfy
\[
 (UV)^{-1}(0)=\{0\},\qquad UV(0)=0,
\]
then, using both factorizations \(UV,VU\),
\[
\begin{alignedat}{2}
 U^{-1}(0)&=\{V(0)\},\qquad&V^{-1}(0)&=\{U(0)\},\\
 U^{-1}(U(0))&=\{0\},&V^{-1}(V(0))&=\{0\}.
\end{alignedat}
\]
For an onto interval map, a value having an ambient endpoint as its unique
preimage must itself be an ambient endpoint; otherwise preimages of the two
range endpoints and the intermediate value theorem give a second preimage.
Apply this to the two factorizations
\[
 S=P\bigl(P^{2L-1}Q^{2L}\bigr)
  =Q\bigl(P^{2L}Q^{2L-1}\bigr),
\]
whose factors are commuting and onto.  It follows that each of \(P,Q\)
either fixes zero or interchanges the ambient endpoints.  Since \(L\) is even,
\[
 A(0)=0=B(0).
\]
Commutation preserves the basin, so \(A(C_e),B(C_e)\subseteq C_e\).  As
\(S=A^2B^2\) and \(S(C_e)=C_e\), the inclusions sandwich to
\[
 A(C_e)=C_e=B(C_e).
\]

The phase images \(P^iQ^j(C_e)\), \(0\leq i,j<L\), now obey the same
vertex-transitive intersection-graph dichotomy.  If their graph is
connected, interval Helly gives a common invariant intersection contained in
the proper interval \(C_e\).  If it is disconnected, collapse the connected
components of their union.  Proposition~\ref{prop:quotient-package}(6) says
that the descended pair is still a counterexample.  Every phase interval
contains the two distinct
images of the endpoints of \(C_e\), which lie in \(\Fix(S)\); distinctness
is preserved because \(P,Q\) have explicit inverses on \(\Fix(S)\).
Thus the quotient strictly lowers \(\#\Fix(S)\), unless an earlier fixed
count has already fallen.  Either outcome contradicts the lexicographic
minimality in \eqref{eq:analytic-mintriple}.  The reflected argument treats
the endpoint one.  Hence no ambient endpoint is merely asymptotic, and
\begin{equation}\label{eq:analytic-finite-arrival}
 R^N(J)=\I
\end{equation}
for some \(N\).

There remains a terminal-gap issue.  Suppose \(c=1\) and \(p>0\).
Equation~\eqref{eq:analytic-gap-sign} makes one an attracting endpoint for
\(R=W^{L/2}\), possibly parabolically.  Indeed, every orbit beginning in
\((p,1)\) increases and converges to the only possible \(W\)-fixed limit,
namely one.  Thus the immediate \(R\)-basin at one has closure
\(C=[p,1]\), which is proper; its free boundary \(p\) is fixed, and
continuity at the two fixed endpoints gives \(R(C)=C\).
Since \(A,B\) fix one and preserve the basin,
\[
 A(C),B(C)\subseteq C,\qquad
 C=AB(C)\subseteq A(C)\subseteq C,
\]
and symmetrically \(B(C)=C\).  The phase-graph dichotomy now gives either a
proper common invariant interval or a quotient strictly lowering
\(\#\Fix(R)\), because every phase interval contains the two distinct
\(R\)-fixed endpoints of \(C\).  In the quotient alternative,
Proposition~\ref{prop:quotient-package}(6) preserves the counterexample.
Both alternatives are impossible.

If \(p=0,c=1\), then consecutiveness of \(p,c\) in \(X\) gives
\(X=\{0,1\}\).  By \eqref{eq:analytic-X-inverses}, the restrictions \(P|_X\) and
\(Q|_X\) are permutations; each therefore either fixes both endpoints or
interchanges them.  Suppose a factor, say \(P\), interchanges the endpoints.
Lemma~\ref{lem:analytic-states} makes its selected germ a local
diffeomorphism from the right of zero to the left of one, hence locally
decreasing.  For small \(x>0\),
\[
 W(x)>x\quad\Longrightarrow\quad P(W(x))<P(x).
\]
Commutation gives \(W(P(x))<P(x)\), contradicting
\eqref{eq:analytic-gap-sign}.  Hence neither factor interchanges the
endpoints; both fix zero, already a common fixed point.  In the surviving
counterexample,
\begin{equation}\label{eq:analytic-c-internal}
 c<1.
\end{equation}

Finally, fix $0<\delta<c-p$, put $x_0=p+\delta$, and write
$x_n=W^n(x_0)$.  As long as $x_n\in(p,c)$,
\eqref{eq:analytic-gap-sign} gives $x_{n+1}>x_n$.  Hence either some
$x_n\geq c$, or the sequence remains in $(p,c)$ and increases to a limit
$d\leq c$.  In the latter case continuity gives $W(d)=d$; since there is no
$W$-fixed point in $(p,c)$, necessarily $d=c$.

Choose $z_0,z_1\in J$ with
\[
 R^N(z_0)=0,\qquad R^N(z_1)=1.
\]
Because $c<1$ and $R$ fixes $p,c$, the point $z_1$ lies in $(p,c)$, and so
does $z_0$ unless $p=0$ and $z_0=p$.  By the preceding orbit argument there
is $n$ such that $x_n$ lies to the right of both $z_0$ and $z_1$ (or already
lies at or to the right of $c$).  The set
$W^n([p,p+\delta])$ is an interval containing
$W^n(p)=p$ and $W^n(x_0)=x_n$; it therefore contains both $z_0$ and $z_1$.
Applying $R^N$, a power of $W$, gives an interval containing both ambient
endpoints, hence all of $\I$.  If $\delta\geq c-p$, the initial interval
contains a smaller ray to which the preceding argument applies.  We have
therefore proved the fundamental local property
\begin{equation}\label{eq:analytic-LEO}
 \boxed{\text{For every }0<\delta\leq1-p\text{ there is }k\text{ such that }
 W^k([p,p+\delta])=\I.}
\end{equation}

In particular, if two maps commuting with \(W\) agree on one selected-side
interval, then they agree globally.  We call this \emph{germ
determination}.

\section{The one-sided \texorpdfstring{\(C^{1+\mathrm{BV}}\)}{C1+BV} orbit circle}

For an orientation-preserving \(C^1\) local diffeomorphism, write
\(C^{1+\mathrm{BV}}\) when \(\log Df\) has bounded variation on a sufficiently
small closed one-sided interval.  This class is closed under composition and
inverse, because
\[
 \log D(g\circ h)=(\log Dg)\circ h+\log Dh,
 \qquad
 \log D(g^{-1})=-(\log Dg)\circ g^{-1}.
\]
Lemma~\ref{lem:analytic-finiteness} shows that every selected germ occurring
here has this regularity: after shrinking the germ domain its positive
derivative is bounded away from zero, so bounded variation of the derivative
implies bounded variation of its logarithm.

We use three standard one-dimensional facts.  Navas's
bounded-variation Kopell lemma \cite[Theorem 4.1.1]{Navas} says that if
commuting orientation-preserving diffeomorphisms $u,v$ of $[0,1)$ are
respectively $C^{1+\mathrm{BV}}$ and $C^1$, if $u$ has no interior fixed
point, and if $v$ has one, then $v=\id$.  The Poincar\'e translation-number
criterion \cite[Propositions 2.2.2 and 2.2.3]{Navas} says that an
orientation-preserving circle-homeomorphism lift $G$ has translation number
$m/n\in\mathbb Q$, $n>0$, exactly when $G^n(t)=t+m$ for some $t$.
Finally, Denjoy's theorem
\cite[Theorem 3.1.1]{Navas} says that an orientation-preserving $C^1$ circle
diffeomorphism whose derivative has bounded variation and whose rotation
number is irrational is topologically conjugate to the corresponding rigid
rotation.

We first record the local form of Kopell's lemma that will be used below.

\begin{lemma}[Local BV Kopell corollary]\label{lem:local-BV-Kopell}
Let \(u,v\) be commuting orientation-preserving one-sided \(C^1\) local
diffeomorphism germs fixing zero.  Suppose that \(u\) is
\(C^{1+\mathrm{BV}}\), has no fixed point on the punctured side, and that
\(v\) has a punctured-side fixed point.  Then \(v\) is the identity germ.
\end{lemma}

\begin{proof}
Replace \(u\) by \(u^{-1}\), if necessary, so that \(u(x)<x\).  Choose a
fixed point \(y>0\) of \(v\) in a common representative domain.  Then every
\(u^j(y)\), \(j\geq0\), is fixed by \(v\), and \(u^j(y)\downarrow0\).
Extend \(u\) above \(y\) to a fixed-point-free
\(C^{1+\mathrm{BV}}\) diffeomorphism of \([0,1)\), preserving the already
given germ.  For completeness, write \(z=u(y)<y\) and \(r=Du(y)>0\).
We construct the extension explicitly.  Put \(a=y-z>0\),
\(s=1-r<1\), and choose
\(0<\eta<\min\{r,1\}/2\).  Fix \(\kappa>0\), and choose
\(Y\in(y,1)\) so close to \(y\) that the function
\[
 q(x)=s+\frac{x-y}{Y-y}(-\kappa-s),\qquad
 d(x)=a+\int_y^x q(t)\,dt \quad (y\leq x\leq Y),
\]
is positive.  Since both endpoint values of \(q\) are at most
\(1-\eta\), so is \(q\) throughout this interval.  Write
\(a_Y=d(Y)>0\) and set
\(\lambda=\kappa(1-Y)^2/a_Y\).  On \([Y,1)\), continue \(d\) by
\[
 d(x)=a_Y\exp\!\left[-\lambda
   \left(\frac1{1-x}-\frac1{1-Y}\right)\right].
\]
The value and first derivative match at \(Y\), because the right derivative
there is \(-\lambda a_Y/(1-Y)^2=-\kappa\).  For the variation calculation,
put \(w=(1-x)^{-1}\).  On the tail,
\[
 d'=-\lambda w^2d,\qquad
 d''=d(\lambda^2w^4-2\lambda w^3).
\]
Since \(d=C e^{-\lambda w}\) there, where
\(C=a_Ye^{\lambda/(1-Y)}\), and \(dx=dw/w^2\),
\[
 \int_Y^1|d''(x)|\,dx
 \leq C\int_{(1-Y)^{-1}}^\infty
 e^{-\lambda w}(\lambda^2w^2+2\lambda w)\,dw<\infty.
\]
Thus \(d'\) has finite variation on the tail.  It equals the affine
function \(q\) on \([y,Y]\), and there is no jump at \(Y\), so \(d'\) has
bounded variation on all of \([y,1)\).  The same formula shows that the tail
and all its derivatives tend to zero at one.  Consequently \(d\) is positive
and \(C^1\), \(d'\leq1-\eta\), and \(d(x)\to0\) as \(x\uparrow1\).  Define
\(u(x)=x-d(x)\) for \(x\geq y\).
Its value and derivative at \(y\) are \(z\) and \(r\), so it glues to the
given representative.  Moreover \(Du=1-d'\geq\eta\), \(u(x)<x\), and
\(u(x)\to1\).  Thus the result is an orientation-preserving diffeomorphism
of \([0,1)\) onto itself with no interior fixed point.  Since \(Du\) has
bounded variation and is bounded away from zero, \(\log Du\) has bounded
variation as well.  This is the required extension.

Keep \(v\) on \([0,y]\).  On the successive fundamental intervals above
\(y\), extend it by
\[
 v\bigm|_{u^{-k}([u(y),y])}
   =u^{-k}\bigl(v\bigm|_{[u(y),y]}\bigr)u^k,
 \qquad k=1,2,\ldots .
\]
Here \(a_k=u^{-k}(y)\) increases to one, and the \(k\)-th displayed formula
acts on \([a_{k-1},a_k]\).  It fixes both endpoints, so the pieces fit
continuously and define an increasing homeomorphism of \([0,1)\).  Indeed,
the increasing sequence \((a_k)\) can have no limit below one, since such a
limit would be an interior fixed point of \(u\).
At the common endpoint \(a_k\), the derivative from the left is \(Dv(y)\),
whereas the derivative from the right is \(Dv(u(y))\): the derivatives of
\(u^{-k}\) and \(u^k\) cancel at their corresponding orbit points.  Since
\(v(y)=y\), differentiating \(uv=vu\) at \(y\) gives
\[
 Du(y)Dv(y)=Dv(u(y))Du(y),
\]
and \(Du(y)>0\), so these two one-sided derivatives agree.  Thus the pieces
glue \(C^1\).  The conjugation formulas also give \(uv=vu\) on each
fundamental interval, and continuity gives the identity at every seam.
Only finitely many seams meet a compact subinterval of \([0,1)\), so the
extension is a \(C^1\) diffeomorphism there and still fixes \(y\).
The BV Kopell lemma makes it the identity, so the original \(v\) is the
identity germ.
\end{proof}

\begin{lemma}[Orbit-circle dichotomy]\label{lem:orbit-circle}
Let \(f,A,B\) be pairwise commuting orientation-preserving
\(C^{1+\mathrm{BV}}\) one-sided local diffeomorphism germs fixing zero, and
suppose
\[
 f(x)>x\quad(0<x\ll1),\qquad AB=f^d,\qquad d\in\mathbb Z_{>0}.
\]
The physical germ \(A\) selects a real lift translation number \(\alpha\),
normalized by \(\tau_f(f)=1\).  Exactly one of the following holds.
\begin{enumerate}
\item If \(\alpha=m/n\in\mathbb Q\), where \(n>0\), then
\begin{equation}\label{eq:rational-germ-powers}
 A^n=f^m,\qquad B^n=f^{dn-m}
\end{equation}
as germs.
\item If \(\alpha\notin\mathbb Q\), there are an increasing homeomorphism
\(\chi:(0,\varepsilon)\to(-\infty,T)\) and simultaneous germ identities,
valid on a sufficiently far-left tail,
\begin{equation}\label{eq:irrational-germ-translations}
 \chi f\chi^{-1}(s)=s+1,\quad
 \chi A\chi^{-1}(s)=s+\alpha,\quad
 \chi B\chi^{-1}(s)=s+d-\alpha.
\end{equation}
\end{enumerate}
\end{lemma}

\begin{proof}
\begin{claim}[Direct-limit orbit line]\label{clm:direct-limit-orbit-line}
For the germs \(f,A,B\) above there is a Hausdorff ordered line
\(\mathscr L\), carrying a \(C^{1+\mathrm{BV}}\) atlas, in which the physical
punctured side is a neighborhood of the negative end.  The three germs
extend uniquely by equivariance to pairwise commuting
\(C^{1+\mathrm{BV}}\) orientation-preserving diffeomorphisms of
\(\mathscr L\), and all their germ relations remain valid.  The extension
of \(f\) acts freely, \(\mathscr L/\langle f\rangle\) is a
\(C^{1+\mathrm{BV}}\) circle, and there is a
\(C^{1+\mathrm{BV}}\) coordinate \(H:\mathscr L\to\mathbb R\) satisfying
\(HfH^{-1}(s)=s+1\).
\end{claim}

\begin{proof}[Proof of the claim]
Shrink the representatives so that $f,A,B$ and their inverses are defined on
the intervals used below.  Here is an explicit model for the direct limit.
For a sufficiently small representative interval \(U=(0,\varepsilon)\),
take classes \([x,j]\), \(x\in U\), \(j\in\mathbb Z\), under the relations
\begin{equation}\label{eq:direct-limit-relation}
 [f(x),j]=[x,j+1]
\end{equation}
whenever the displayed terms are defined.  The ordered direct limit is an
open line \(\mathscr L\).  The physical punctured side embeds by
\(x\mapsto[x,0]\); because \(f^{-k}(x)\downarrow0\), it is a neighborhood of
the negative end.  Using a sufficiently small representative of each class,
the germ extends to the free deck map

\[
 f_{\mathscr L}[x,j]=[f(x),j]=[x,j+1].
\]

We continue to denote this extension by \(f\).  Choose \(x_0>0\) and put
\(x_1=f(x_0)>x_0\).  Every point of \(\mathscr L\) has a representative
\([x,j]\) with \(x\in D=[x_0,x_1]\); the only boundary identifications are
\([x_1,j]=[x_0,j+1]\).  Thus \(\mathscr L\) is the ordered gluing of the
copies \(D_j=D\times\{j\}\).  Reduction to \(D\) is unique except for these
two seam representatives, and
\[
 H_0([x,j])=j+\frac{x-x_0}{x_1-x_0}\qquad(x\in D)
\]
is well defined at each seam.  Its restriction to every copy is a
homeomorphism onto \([j,j+1]\), and the piecewise inverse has the quotient
topology.  Hence \(H_0\) is a homeomorphism
\(\mathscr L\cong\mathbb R\); in particular the quotient is Hausdorff and
the asserted ordered open line is genuine.

For regularity, enlarge each copy slightly across its two endpoints using
\eqref{eq:direct-limit-relation}.  The overlap maps are restrictions of
iterates of \(f\) and \(f^{-1}\), so they form a
\(C^{1+\mathrm{BV}}\) atlas compatible with this topology.  The deck map
is represented in these charts by \(f\) and its conjugates.  Its quotient
\[
 \mathscr O_f=\mathscr L/\langle f\rangle
\]
is a topological circle.  A fundamental interval is
$D=[x_0,x_1]$, with $x_0$ identified with $x_1$ by $f$.

We spell out the regularity at this seam.  Choose a positive $C^1$ function
$\rho$ on $D$, with derivative of bounded variation, satisfying
\[
 \rho(x_1)Df(x_0)=\rho(x_0).
\]
For example, prescribe the two positive endpoint values in this ratio and
join them by a positive affine function.  Rescale \(\rho\) so that
\(\int_{x_0}^{x_1}\rho=1\), and
set
\[
 h(x)=\int_{x_0}^x\rho(t)\,dt,\qquad x\in D.
\]
Extend $h$ to a coordinate $H:\mathscr L\to\mathbb R$ by

\[
 H([x,j])=h(x)+j \qquad (x\in D).
\]

The values agree at the seam because \(h(x_1)=1\) and \(h(x_0)=0\).
To check the derivative, use the \(j\)-th physical coordinate on both sides.
Just to the right of the seam, write the point as
\([x,j+1]=[f(x),j]\).  The derivative there, in the coordinate \(f(x)\),
tends to \(\rho(x_0)/Df(x_0)\); from the left it tends to \(\rho(x_1)\).
These are equal by the endpoint condition on \(\rho\).  Hence there is no
derivative jump.  Each side is \(C^{1+\mathrm{BV}}\), and a compact part of
\(\mathscr L\) meets only finitely many copies, so the total variation on
that compact part is finite.  Positivity of \(\rho\) gives the same
statement for the local inverse.  Thus \(H\) is a
\(C^{1+\mathrm{BV}}\) coordinate and
\begin{equation}\label{eq:orbit-deck}
 HfH^{-1}(s)=s+1.
\end{equation}
It follows in particular that $\mathscr O_f$ carries the standard
$C^{1+\mathrm{BV}}$ circle structure.

For \(G=A\) or \(G=B\), define its equivariant extension by
\begin{equation}\label{eq:direct-limit-G}
 G_{\mathscr L}[x,j]=[G(x),j],
\end{equation}
using a sufficiently small representative of the class.  This is
independent of the representative: relation
\eqref{eq:direct-limit-relation} and \(Gf=fG\) give

\[
 [G(f(x)),j]=[f(G(x)),j]=[G(x),j+1].
\]

The same construction with the inverse germ proves that
\(G_{\mathscr L}\) is an orientation-preserving diffeomorphism, and the
definition gives \(G_{\mathscr L}f=fG_{\mathscr L}\).  In overlapping
direct-limit charts its formulas are conjugates \(f^kGf^{-k}\).  Adjacent
formulas have equal values by commutation, and differentiating \(Gf=fG\)
shows that their first derivatives agree at the seam.  All such formulas
are \(C^{1+\mathrm{BV}}\).  Since a compact interval in \(\mathscr L\)
meets only finitely many charts, its derivative variation is a finite sum
of the piecewise variations, with no seam jump.  Thus \(G_{\mathscr L}\)
is \(C^{1+\mathrm{BV}}\) on the line.
Formula \eqref{eq:direct-limit-G} also preserves every germ relation:
evaluate both words on a representative sufficiently close to zero.  In
particular the extensions of \(A,B,f\) commute pairwise and satisfy
\(AB=f^d\).  Equivariance and agreement on the embedded negative end make
the extensions unique.  This proves the claim.
\end{proof}

Choose the line, extensions, and coordinate supplied by
Claim~\ref{clm:direct-limit-orbit-line}.  For \(G\in\{A,B\}\), put
\(\widetilde G=HG_{\mathscr L}H^{-1}\).  Commutation with \(f\) gives the
exact deck identity
\[
 \widetilde G(s+1)=\widetilde G(s)+1.
\]
Consequently \(D\widetilde G(s+1)=D\widetilde G(s)\); restricting to
\([0,1]\) and identifying its endpoints produces an orientation-preserving
\(C^{1+\mathrm{BV}}\) circle diffeomorphism.  This proves the required
regularity of the circle maps induced by \(A\) and \(B\).

The actual action of $G$ on $\mathscr L$, rather than only its quotient
circle map, selects this lift $\widetilde G$ commuting with the deck
translation $T_1(s)=s+1$.  A different lift $T_1^k\widetilde G$ would
represent the different germ $f^kG$.  Define
\[
 \alpha=\tau_f(A)=\lim_{r\to\infty}
       \frac{\widetilde A^{\,r}(s)-s}{r}.
\]
This is independent of $s$, its class modulo $\mathbb Z$ is the rotation
number of the induced circle map, and the germ identity $AB=f^d$ gives the
exact lift identity
\begin{equation}\label{eq:selected-lift-product}
 \widetilde A\widetilde B=T_1^d.
\end{equation}

Suppose first that $\alpha=m/n\in\mathbb Q$, with $n>0$.  The
Poincar\'e theorem provides $s$ such that
$\widetilde A^n(s)=s+m$.  Translating $s$ by a sufficiently negative integer
places the corresponding point in the physical representative of the germ;
the equality is unchanged because the lifts commute with $T_1$.  Hence
$H_0=A^nf^{-m}$ has a punctured-side fixed point.  It commutes with the
fixed-point-free $C^{1+\mathrm{BV}}$ germ $f$, so
Lemma~\ref{lem:local-BV-Kopell} gives $H_0=\id$.  Therefore
$A^n=f^m$.  Since $B=f^dA^{-1}$, also
$B^n=f^{dn-m}$.

Now suppose that $\alpha$ is irrational.  Denjoy's theorem conjugates the
circle map induced by $A$ to rotation by $\alpha\bmod\mathbb Z$.  Lift a
degree-one conjugacy to an increasing homeomorphism of universal covers
which commutes with $T_1$.  The conjugated selected lift of $A$ is a
translation.  Its displacement is exactly $\alpha$, rather than merely its
residue class, because translation number is invariant under degree-one lift
conjugacy.  Equation~\eqref{eq:selected-lift-product} then sends the selected
lift of $B$ to translation by $d-\alpha$, while
\eqref{eq:orbit-deck} sends $f$ to translation by one.  Restricting this
lifted coordinate to a sufficiently far-left part of the physical end gives
an increasing homeomorphism
$\chi:(0,\varepsilon)\to(-\infty,T)$ and the three identities in
\eqref{eq:irrational-germ-translations}.
\end{proof}

Apply Lemma~\ref{lem:orbit-circle} in the physical selected-side coordinate
\(x-p\) to \(f=W\) and to \(A,B\), where \(d=L/2\) and \(AB=W^d\).
There is also an elementary one-map dilation chart
\[
 \jmath:[0,\varepsilon_0)\longrightarrow[p,p+\delta_0),\qquad
 W(\jmath(t))=\jmath(et)\quad(0\leq t<\varepsilon_0/e).
\]
To construct it, choose an arbitrary increasing homeomorphism between one
pair of fundamental intervals and extend it equivariantly by the iterates of
\(W\), then set \(\jmath(0)=p\).

Suppose first that \(\alpha=m/n\) is rational.  If \(m<0\), choose a small
\(0<\varepsilon<\varepsilon_0/e\) and put
\(I_\varepsilon=\jmath([0,\varepsilon])\), shrinking \(\varepsilon\) until
the germ relation holds there.
Here a negative power of \(W\) denotes the corresponding iterate of the
selected local inverse germ; no global inverse of \(W\) is being used.
Then
\[
 A^{nk}(I_\varepsilon)=W^{mk}(I_\varepsilon)\longrightarrow\{p\}.
\]
Choose \(q\) from \eqref{eq:analytic-LEO} with
\(W^q(I_\varepsilon)=\I\).  Onto-ness and commutation would give
\[
 \I=A^{nk}(\I)=W^q(A^{nk}(I_\varepsilon))\longrightarrow\{p\},
\]
a contradiction.  If \(m=0\), germ determination gives \(A^n=\id\)
globally, so \(P^{Ln}=\id\).  Then \(P\) is a finite-order interval
homeomorphism: if it preserves orientation it is the identity, while if it
reverses orientation its unique fixed point is preserved by \(Q\).  Either
case gives a common fixed point.  Applying the same argument to
\(B^n=W^{dn-m}\) shows that \(dn-m\) is neither negative nor zero.  Hence
\(0<m<dn\), and
\[
 A^{\,n(dn-m)}=B^{\,nm}.
\]
Germ determination makes this equality global, that is,
\(P^{Ln(dn-m)}=Q^{Lnm}\).  Both powers are positive, so
Lemma~\ref{lem:analytic-common-iterate} below again gives a common fixed
point of \(P,Q\).  Thus the rational alternative is impossible in the
standing counterexample.

We are in the irrational alternative.  Exponentiate its translation
coordinate and return to the physical side.  This gives the germ chart
\[
 \iota:[0,e^T)\longrightarrow\I,\qquad
 \iota(0)=p,\qquad
 \iota(t)=p+\chi^{-1}(\log t)\quad(0<t<e^T),
\]
and set
\[
 \lambda=e,\qquad a=e^\alpha,qquad b=e^{d-\alpha},qquad
 \Lambda=e^d=\lambda^d.
\]
Then, as selected-side germ identities (that is, whenever \(t\) and the
scaled arguments displayed below lie in \([0,e^T)\)),
\begin{equation}\label{eq:actor-dilations}
 W(\iota(t))=\iota(\lambda t),\qquad
 A(\iota(t))=\iota(at),\qquad
 B(\iota(t))=\iota(bt),
\end{equation}
and
\begin{equation}\label{eq:analytic-scale-product}
 ab=\Lambda>1.
\end{equation}

No onto centralizer return can have scale below one.  If, for example,
\(0<a<1\), then for a small chart interval
\(I_\varepsilon=\iota([0,\varepsilon])\),
\(A^k(I_\varepsilon)=\iota([0,a^k\varepsilon])\to\{p\}\).  Choosing \(q\)
with \(W^q(I_\varepsilon)=\I\), surjectivity and commutation give
\(\I=A^k(\I)=W^q(A^k(I_\varepsilon))\to\{p\}\), a contradiction.  Here
\(a=1\) would make \(\alpha=0\), and \(b=1\) would make \(\alpha=d\); both
contradict irrationality.  Thus
\begin{equation}\label{eq:analytic-scales-positive}
 a,b>1.
\end{equation}

We shall use the Fugate--McLean fixed-set theorem
\cite[Corollary 1.5, p.~614]{FugateMcLean}, whose exact conclusion is that
the fixed-point set of a periodic homeomorphism of a tree-like continuum is
a nonvoid continuum.  Thus, in particular, \(\Fix(V)\) is nonempty, compact,
and connected.

\begin{lemma}[Common positive iterate]
\label{lem:analytic-common-iterate}
Let \(F,G\colon\I\to\I\) be commuting continuous maps.  If
\[
 F^m=G^n
\]
for some positive integers \(m,n\), then \(F,G\) have a common fixed point.
\end{lemma}

\begin{proof}
Put \(T=F^m=G^n\) and
\[
 J_\infty=\bigcap_{r\geq0}T^r(\I).
\]
The sets in this intersection are nested nonempty compact intervals.  Also
\begin{equation}\label{eq:analytic-eventual-range-onto}
 T(J_\infty)=J_\infty.
\end{equation}
Indeed, one inclusion is immediate.  For \(y\in J_\infty\), the nonempty
compact sets
\[
 T^{-1}(y)\cap T^r(\I),\qquad r\geq0,
\]
are nested: nonemptiness follows from \(y\in T^{r+1}(\I)\).  Their
intersection contains a preimage of \(y\) in \(J_\infty\), proving the other
inclusion.  Commutation gives
\[
 F(J_\infty)\subseteq J_\infty,\qquad
 G(J_\infty)\subseteq J_\infty.
\]
If \(J_\infty\) is a singleton, it is already common fixed.  Assume it is
nondegenerate.

Form the natural extension
\[
 X_T=\{(x_0,x_1,\ldots)\in J_\infty^{\mathbb N_0}:
       T(x_{i+1})=x_i\text{ for every }i\geq0\}.
\]
By \eqref{eq:analytic-eventual-range-onto}, this is a nonempty compact
connected metric space.  Give it the metric
\[
 d_X(x,y)=\sum_{i=0}^{\infty}2^{-i-1}|x_i-y_i|.
\]
The coordinate projection \(\pi_N\colon X_T\to J_\infty\) is onto.  If
\(\pi_N(x)=\pi_N(y)\), the bonding equations make \(x_i=y_i\) for
\(0\leq i\leq N\), and hence
\[
 d_X(x,y)\leq\sum_{i>N}2^{-i-1}=2^{-N-1}.
\]
Thus, for every \(\varepsilon>0\), some \(\pi_N\) is an
\(\varepsilon\)-map onto an interval.  Hence \(X_T\) is arc-like and
therefore tree-like.

The shift
\[
 \Sigma(x_0,x_1,\ldots)=(T(x_0),x_0,x_1,\ldots)
\]
is a homeomorphism, with inverse deleting the zeroth coordinate.
Coordinatewise application defines commuting maps \(\widehat F,\widehat G\)
of \(X_T\), and
\[
 \widehat F^m=\Sigma=\widehat G^n.
\]
A continuous selfmap whose positive power is a homeomorphism is itself a
homeomorphism: injectivity follows from injectivity of the power, surjectivity
from surjectivity of the power, and a continuous bijection of a compact space
onto a Hausdorff space is a homeomorphism.  Thus \(\widehat F,\widehat G\)
are commuting homeomorphisms.

Write
\[
 m=d_0m_0,\qquad n=d_0n_0,\qquad \gcd(m_0,n_0)=1,
\]
and choose integers \(r_0,s_0\) with
\[
 r_0n_0+s_0m_0=1.
\]
Negative powers are now legitimate.  Set
\[
 U=\widehat F^{r_0}\widehat G^{s_0},\qquad
 V=\widehat F^{m_0}\widehat G^{-n_0}.
\]
Then \(V^{d_0}=\id\), and the B\'ezout identity gives
\begin{equation}\label{eq:analytic-bezout-split}
 \widehat F=U^{n_0}V^{s_0},\qquad
 \widehat G=U^{m_0}V^{-r_0}.
\end{equation}

The quoted Fugate--McLean theorem makes
\[
 \mathcal C=\Fix(V)
\]
a nonempty compact connected subset of \(X_T\).  Since \(U\) commutes with
\(V\), it maps \(\mathcal C\) into itself.  We prove directly that
\(U|_{\mathcal C}\) has a fixed point.  Given \(\varepsilon>0\), choose
\(N\) with \(2^{-N-1}<\varepsilon\).  The set
\(\pi_N(\mathcal C)=[\alpha,\beta]\) is an interval.  If
\(\alpha=\beta\), every \(z\in\mathcal C\) satisfies
\(d_X(U(z),z)<\varepsilon\).  Otherwise choose
\(z_\alpha,z_\beta\in\mathcal C\) over the two endpoints.  The continuous
function
\[
 z\longmapsto\pi_N(U(z))-\pi_N(z)
\]
is nonnegative at \(z_\alpha\) and nonpositive at \(z_\beta\).  Its image of
the connected set \(\mathcal C\) therefore contains zero.  In either case
there is \(z_\varepsilon\in\mathcal C\) with
\[
 d_X(U(z_\varepsilon),z_\varepsilon)<\varepsilon.
\]
Compactness, followed by \(\varepsilon\downarrow0\), supplies
\(z\in\mathcal C\) with \(U(z)=z\).  Since also \(V(z)=z\),
\eqref{eq:analytic-bezout-split} gives
\(\widehat F(z)=z=\widehat G(z)\).  The zeroth coordinate is a common fixed
point of \(F,G\).
\end{proof}

If \(a^m=b^n\) for positive integers \(m,n\), the dilation coordinate makes
\(A^m\) and \(B^n\) equal as germs.  Germ determination makes them equal
globally:
\[
 P^{Lm}=Q^{Ln}.
\]
 Lemma~\ref{lem:analytic-common-iterate} then gives a common fixed point of
 \(P,Q\).
Consequently the sole remaining case is
\begin{equation}\label{eq:analytic-independent-scales}
 a,b>1,\qquad \frac{\log a}{\log b}\notin\mathbb Q.
\end{equation}

\section{The canonical ray and its tail core}

For \(t\geq0\), choose \(n\) large enough that \(t/\lambda^n\) lies in the
domain of \(\iota\), and define
\begin{equation}\label{eq:analytic-ray}
 \Phi(t)=W^n\left(\iota\left(\frac{t}{\lambda^n}\right)\right).
\end{equation}
Equation~\eqref{eq:actor-dilations} makes this independent of every
sufficiently large choice of \(n\).  One \(n\) works on each compact
parameter interval, so \(\Phi\) is continuous; near zero it equals
\(\iota\).  Property~\eqref{eq:analytic-LEO} makes \(\Phi\) onto.  Commutation
and the germ identities give
\begin{equation}\label{eq:analytic-ray-equivariance}
 \Phi(at)=A(\Phi(t)),\qquad
 \Phi(bt)=B(\Phi(t)),\qquad
 \Phi(\Lambda t)=R(\Phi(t)).
\end{equation}

For \(T>0\), put
\[
 \mathcal C_T=\cl\bigl(\Phi([T,\infty))\bigr),\qquad
 \Omega=\bigcap_{T>0}\mathcal C_T.
\]
Each \(\mathcal C_T\) is a compact interval, and equivariance gives
\(A(\mathcal C_T)=\mathcal C_{aT}\).  Here is the closure bookkeeping.
If \(x=\lim_k\Phi(t_k)\in\mathcal C_T\), with \(t_k\geq T\), continuity and
equivariance give
\(A(x)=\lim_k\Phi(at_k)\in\mathcal C_{aT}\).  Conversely, if
\(y=\lim_k\Phi(s_k)\in\mathcal C_{aT}\), take \(s_k\geq aT\) and put
\(t_k=s_k/a\).  A subsequence of \(\Phi(t_k)\) converges, by compactness of
\(\I\), to some \(x\in\mathcal C_T\), and then
\(A(x)=\lim_k\Phi(s_k)=y\).  This proves the asserted equality.  Hence
\(A(\Omega)\subseteq\Omega\).  Conversely, if \(y\in\Omega\), then the
nested compact sets
\[
 \mathcal C_T\cap A^{-1}(y)\qquad(T>0)
\]
are nonempty because \(y\in\mathcal C_{aT}=A(\mathcal C_T)\).  Their
intersection contains \(x\in\Omega\) with \(A(x)=y\).  Thus
\(A(\Omega)=\Omega\), and the identical argument for \(B\) gives
\begin{equation}\label{eq:analytic-Omega-invariant}
 A(\Omega)=\Omega=B(\Omega).
\end{equation}

Suppose first that \(\Omega=\{q\}\).  If \(x\) is fixed by \(R\), choose
\(t\) with \(\Phi(t)=x\).  If \(t>0\), then
\[
 x=R^n(x)=\Phi(\Lambda^nt)\longrightarrow q;
\]
the lift \(t=0\) gives \(x=p\).  Hence
\[
 \Fix(R)\subseteq\{p,q\}.
\]
The adjacent point \(c\) is fixed by \(A,B\), so \(q=c\).
The sign of \(R-\id\) is positive on \((p,q)\) and negative on
\((q,1]\): it is positive near \(p\), and there is no further fixed point
on either interval.

Commutation gives \(P(p),Q(p)\in\Fix(R)\), so each is either \(p\) or \(q\).
If \(P(p)=q\), its nonzero selected germ maps the right side of \(p\) to one
of the two sides at \(q\).  On the source, \(R\) moves points away from \(p\).
Commutation would therefore make the target side repelling.  But on the left
of \(q\), \(R(x)>x\) moves points toward \(q\), while on the right,
\(R(x)<x\) does the same.  This order contradiction proves \(P(p)=p\).
Similarly \(Q(p)=p\), contrary to the assumption of no common fixed point.

Suppose instead that \(\Omega\) is nondegenerate and proper.  Its phase
images
\[
 P^iQ^j(\Omega),\qquad0\leq i,j<L,
\]
are nondegenerate: if one were a singleton, applying
\(P^{L-i}Q^{L-j}\) would make
\(P^LQ^L(\Omega)=R(\Omega)=\Omega\) a singleton.  They have a
vertex-transitive intersection graph by
\eqref{eq:analytic-Omega-invariant} and the phase-intersection argument
above.  If the
graph is connected, interval Helly gives a proper common invariant
intersection.  If it is disconnected, collapse the connected components of
the union.  Since \(R\) fixes every phase label, each collapsed component is
mapped onto itself by \(R\).  Every continuous onto selfmap \(f\) of a
nondegenerate interval has at least two fixed points of \(f^2\).  Indeed, if
\(f\) has at least two fixed points there is nothing to prove.  Otherwise
let \(z\) be its unique fixed point.  It is interior, and
\(f(x)>x\) to its left while \(f(x)<x\) to its right.  If the interval is
\([a_0,b_0]\), surjectivity gives \(v<z<u\) with
\(f(v)=b_0\) and \(f(u)=a_0\).  Choose \(x_0\in[v,z)\) with \(f(x_0)=u\).
Then
\[
 f^2(x_0)-x_0=a_0-x_0\leq0,\qquad
 f^2(a_0)-a_0\geq0.
\]
Equality, or the intermediate value theorem, gives an \(f^2\)-fixed point
to the left of \(z\).  Thus every collapsed component
contains two points of \(\Fix(S)\), so the quotient strictly lowers the
third count in \eqref{eq:analytic-mintriple}, unless an earlier count has
already fallen.  Proposition~\ref{prop:quotient-package}(6) preserves the
counterexample under this quotient, so this is impossible.

We conclude that
\begin{equation}\label{eq:analytic-Omega-full}
 \Omega=\I.
\end{equation}
Equivalently, every late ray tail has dense image.

\section{Dense tails imply exactness}

Recall that an interval map \(T\) is \emph{topologically exact} if every
nonempty relatively open set \(U\subseteq\I\) satisfies
\(T^n(U)=\I\) for some \(n\geq0\).

We first produce a positive boundary lift of the base value \(p\).
If \(p>0\), every late tail image is dense and connected, hence contains
values on both sides of \(p\).  A boundary point \(t_*>0\) of a component of
\(\{t:\Phi(t)>p\}\) satisfies
\begin{equation}\label{eq:analytic-boundary-lift}
 \Phi(t_*)=p
\end{equation}
and is approached by values on the selected right side.

Suppose \(p=0\).  If \(\Phi^{-1}(0)=\{0\}\), then
\(R^{-1}(0)=\{0\}\): for \(R(x)=0\), choose \(x=\Phi(t)\); then
\(\Phi(\Lambda t)=0\), forcing \(t=0\).  Choose a small chart interval
\[
 I_\varepsilon=\iota([0,\varepsilon]).
\]
Choose \(\eta>0\) so that \(\iota\) is defined on
\([0,\Lambda\eta]\) and
\[
 R(\iota(t))=\iota(\Lambda t)\qquad(0\leq t\leq\eta).
\]
Since \(R^{-1}(0)=\{0\}\), compactness gives
\[
 \delta:=\min_{x\in[\iota(\eta),1]}R(x)>0.
\]
Now reduce \(\varepsilon\) so that
\(\varepsilon<\Lambda\eta\) and \(\iota(\varepsilon)<\delta\).
On \([0,\iota(\eta)]\) the displayed dilation identifies the full preimage
of \(I_\zeta\), while the compact-complement bound excludes every other
preimage.  Consequently, for \(0<\zeta\leq\varepsilon\),
\begin{equation}\label{eq:analytic-local-inverse}
  R^{-1}(I_\zeta)=I_{\zeta/\Lambda}.
\end{equation}
If \(\Phi(t)\in I_\varepsilon\), then
\[
 \Phi(t)=R^N\bigl(\Phi(t/\Lambda^N)\bigr).
\]
Iterating \eqref{eq:analytic-local-inverse} gives
\[
 \Phi(t/\Lambda^N)\in I_{\varepsilon/\Lambda^N}.
\]
For large \(N\), the ray equals the initial chart at \(t/\Lambda^N\), so
strict monotonicity of \(\iota\) yields \(t\leq\varepsilon\).  Conversely,
\(\Phi=\iota\) on a sufficiently small initial interval.  After reducing
\(\varepsilon\), therefore,
\[
 \Phi^{-1}(I_\varepsilon)=[0,\varepsilon],
\]
contradicting density of every late tail.  Thus the zero fibre has a positive
point.  It cannot contain an entire tail, so a boundary point of the fibre
gives \(t_*>0\) with the same selected-side property as above.

Equivariance and the positive selected germs show that every
\begin{equation}\label{eq:analytic-actor-boundaries}
 a^rb^s t_*,\qquad r,s\in\mathbb Z_{\geq0},
\end{equation}
has the same boundary property.

We now prove exactness of \(R\).  Let \(U\) be nonempty and relatively open.
Tail density gives a compact parameter interval
\[
 J_0=[u,v],\qquad0<u<v,\qquad\Phi(J_0)\subseteq U.
\]
By \eqref{eq:analytic-independent-scales}, the multiplicative group
\[
 \{a^rb^s:r,s\in\mathbb Z\}
\]
is dense in \((0,\infty)\).  Choose integers \(r,s\) such that
\[
 z=t_*a^rb^s\in(u,v).
\]
For \(k\geq\max\{0,-r,-s\}\), the point
\[
 \Lambda^kz=t_*a^{r+k}b^{s+k}
\]
is a selected-side boundary lift in the interior of \(\Lambda^kJ_0\).
Consequently
\[
 R^k(\Phi(J_0))=\Phi(\Lambda^kJ_0)
\]
contains a nondegenerate selected-side interval based at \(p\).  A further
power of \(R=W^d\) maps that interval onto \(\I\): round the exponent in
\eqref{eq:analytic-LEO} upward to a multiple of \(d\), using surjectivity of
\(W\).  Hence \(R\) is topologically exact.

We prove exactness of \(A\) separately.  Given a nonempty relatively open
\(U\), tail density supplies an arbitrarily late compact interval
\[
 J=[u,v],\qquad t_*<u<v,\qquad\Phi(J)\subseteq U.
\]
Let \(j\) be least with \(z=\Lambda^jt_*>u\).  Then
\[
 u<z\leq\Lambda u<\Lambda v.
\]
The set of \(\rho\in(1,\Lambda)\) for which
\[
 \rho u<z<\rho v
\]
contains a nonempty compact interval \(K_0\) in its interior.  For every
\(\rho\in K_0\), the interval \(\rho J\) contains the boundary lift \(z\)
in its interior.  Thus \(\Phi(\rho J)\) contains a nondegenerate selected
interval based at \(p\).  To make the uniformity explicit, define
\[
 m(\rho)=\max\{\Phi(t):t\in[\rho u,\rho v]\}.
\]
The change of variables \(t=\rho((1-s)u+sv)\) puts this maximum over the
fixed compact set \(s\in[0,1]\).  Uniform continuity of \(\Phi\) on the
compact union of these intervals therefore makes \(m\) continuous on
\(K_0\).  The boundary property at \(z\), together with
\(z\in(\rho u,\rho v)\), gives \(m(\rho)>p\) for every \(\rho\in K_0\).
Consequently
\(\eta_0=\min_{\rho\in K_0}(m(\rho)-p)>0\).  Since
\(\Phi([\rho u,\rho v])\) is an interval containing both \(p\) and
\(m(\rho)\), we have
\begin{equation}\label{eq:analytic-uniform-selected}
 [p,p+\eta_0]\subseteq\Phi(\rho J)
 \qquad(\rho\in K_0).
\end{equation}
Choose \(q\) with \(R^q([p,p+\eta_0])=\I\).  Since \(R\) is onto, the same
holds for every exponent at least \(q\).

Write uniquely
\[
 a^n=\Lambda^{k_n}\rho_n,\qquad1\leq\rho_n<\Lambda.
\]
The residues \(\rho_n\) are dense in the multiplicative circle
\([1,\Lambda]\), because
\[
 \frac{\log a}{\log\Lambda}
 =\frac{\log a}{\log a+\log b}
\]
is irrational, and \(k_n\to\infty\).  Choose \(n\) with
\(\rho_n\in K_0\) and \(k_n\geq q\).  Then
\[
\begin{split}
 A^n(\Phi(J))
 &=\Phi(a^nJ)
 =R^{k_n}(\Phi(\rho_nJ))\\
 &\supseteq R^{k_n}([p,p+\eta_0])=\I.
\end{split}
\]
Since \(\Phi(J)\subseteq U\), this proves exactness of \(A\).  Replacing
\(a\) by \(b\) proves exactness of \(B\).  Finally \(A=P^L\) and \(B=Q^L\),
so \(P,Q\) are exact as well.  In particular, none of
\(P,Q,A,B,R\) has a flat interval.

\section{A common Parry coordinate and the polygonal completion}

By Lemma~\ref{lem:analytic-finiteness}, the flat-free tame maps
\(P,Q,R\) have finitely many strict monotonicity laps and are finite-to-one.
We use the following combination of Parry's conjugacy theorem
\cite{Parry} and Shultz's published scaled-measure results
\cite[Corollaries 4.4 and 4.6 and Proposition 3.6]{Shultz}.  A
piecewise-monotone interval map is
\emph{essentially injective} when there are no two disjoint nondegenerate
intervals with the same image.

\medskip
\noindent\textbf{Parry--Shultz scaled-measure theorem.}
\emph{Let $T$ be a topologically transitive piecewise-monotone map with
finitely many laps on a compact interval.  Then it admits a nonatomic
full-support Borel
probability measure $\mu$ and a number
$s=\exp(h_{\rm top}(T))\geq1$ such that
\[
 \mu(T(E))=s\mu(E)
\]
for every Borel set $E$ on which $T$ is injective.  Its distribution function
conjugates $T$ to a finite-piece piecewise-linear map whose nonzero slopes
have absolute value $s$.  If $T$ is not essentially injective, the scaled
probability is unique.  If $T$ is continuous and transitive, then $s>1$.}
\medskip

This theorem applies to $R$.  Exactness implies transitivity.  An exact
interval map cannot be monotone: an onto flat-free monotone map is a
homeomorphism, and
no iterate of a sufficiently small open interval under a homeomorphism is
the whole interval.  Hence $R$ has a genuine turn.  At a genuine turn the
two adjacent strict laps have subintervals with the same nondegenerate image,
so $R$ is not essentially injective in Shultz's sense.  Let $\mu$ be the
unique scaled probability and let $s>1$ be its scale.

For a finite-to-one map \(T\), define the preimage transfer operator
\[
 (\mathcal L_T\varphi)(y)=\sum_{Tx=y}\varphi(x).
\]
Regrouping finite fibres gives
\begin{equation}\label{eq:analytic-transfer-composition}
 \mathcal L_{U\circ T}=\mathcal L_U\mathcal L_T.
\end{equation}

\begin{lemma}[Transfer to a centralizer]\label{lem:analytic-transfer}
If \(T\) is an onto, flat-free, piecewise-monotone map with finitely many laps
and commuting with \(R\), then there is a constant \(c_T>0\) such that
\[
 \mathcal L_T^*\mu=c_T\mu.
\]
\end{lemma}

\begin{proof}
First partition \(\I\) into finitely many pairwise disjoint Borel lap sets
\(I_1,\ldots,I_N\), assigning every lap endpoint to one adjacent set, so
that \(R\) is injective on each \(I_j\).  Since \(\mu\) is nonatomic, the
finite endpoint assignment is immaterial.  For every Borel set \(E\),
\[
 \mathcal L_R\mathbf1_E
 =\sum_{j=1}^N\mathbf1_{R(E\cap I_j)}.
\]
The scaled-set identity therefore gives
\[
\begin{split}
 (\mathcal L_R^*\mu)(E)
 &=\int\mathcal L_R\mathbf1_E\,d\mu\\
 &=\sum_{j=1}^N\mu(R(E\cap I_j))
 =s\sum_{j=1}^N\mu(E\cap I_j)
 =s\mu(E).
\end{split}
\]
Conversely, this eigenmeasure identity implies the scaled-set identity on
every injectivity set.  Thus scaling is equivalent to
\begin{equation}\label{eq:analytic-R-eigenmeasure}
 \mathcal L_R^*\mu=s\mu.
\end{equation}

Put \(\nu_T=\mathcal L_T^*\mu\).  If \(T\) has \(M\) laps, then
\(\mathcal L_T\mathbf1\leq M\), so \(\nu_T\) is finite.  Onto-ness gives
\(\mathcal L_T\mathbf1\geq1\), so it is nonzero.  It is nonatomic because
\[
 \nu_T(\{x\})=\mu(\{T(x)\})=0.
\]
It has full support: every nonempty open interval contains a smaller interval
\(J\) on which \(T\) is strictly monotone, and
\[
 \nu_T(J)=\mu(T(J))>0.
\]

Commutation and \eqref{eq:analytic-transfer-composition} yield
\[
 \mathcal L_R^*\nu_T
 =\mathcal L_T^*\mathcal L_R^*\mu
 =s\nu_T.
\]
After division by its positive total mass, \(\nu_T\) is another nonatomic
full-support scaled probability for \(R\) with factor \(s\).  Uniqueness
gives \(\nu_T=c_T\mu\).
\end{proof}

Apply Lemma~\ref{lem:analytic-transfer} to \(T=P,Q\), and write
\[
 \mathcal L_P^*\mu=C\mu,\qquad
 \mathcal L_Q^*\mu=D\mu
\]
with \(C,D>0\).  Define the distribution coordinate
\begin{equation}\label{eq:analytic-distribution}
 \theta(x)=\mu([0,x]).
\end{equation}
Full support and atomlessness make \(\theta\) an increasing homeomorphism of
\(\I\).

If \(J=[x,y]\) lies in a strict monotonicity lap of \(P\), the transfer
identity gives
\[
 \mu(P(J))=C\mu(J),
\]
and hence
\[
 |\theta(P(y))-\theta(P(x))|
 =C|\theta(y)-\theta(x)|.
\]
It follows that \(\theta P\theta^{-1}\) is affine with slope \(+C\) or
\(-C\) on each of finitely many laps.  Thus it is a finite-piece PL map.
The same argument makes \(\theta Q\theta^{-1}\) finite-piece PL.

These two conjugated maps commute.

Put
\[
 p_\theta=\theta(p).
\]
From now on we use the same letters \(P,Q,A,B,R\) for the maps conjugated by
\(\theta\).  Exactness and commutation are preserved by conjugacy.  Thus
\(P,Q,A,B,R\) are exact, flat-free, finite-piece piecewise-linear maps, with
\[
 A=P^L,\qquad B=Q^L,\qquad R=AB.
\]
Set
\begin{equation}\label{term:eq:transformed-scales}
 \widehat a=C^L,\qquad
 \widehat b=D^L,\qquad
 \widehat s=\widehat a\widehat b.
\end{equation}
Every nonzero slope of \(A,B,R\) has absolute value
\(\widehat a,\widehat b,\widehat s\), respectively.  Since the original
selected germs of \(A,B\) preserve the selected side, their conjugated germs
at \(p_\theta\) have positive slopes \(\widehat a,\widehat b\); consequently
the selected germ of \(R\) has positive slope \(\widehat s\).

\subsection{Multiplicative independence survives the Parry coordinate}

The original selected germs of \(A,B\) move every sufficiently close point
strictly away from \(p\).  Since \(\theta\) preserves order, their conjugated
affine germs move every sufficiently close point strictly away from
\(p_\theta\).  Their positive slopes therefore satisfy
\(\widehat a,\widehat b>1\), and hence \(C,D>1\).

If \(\widehat a^{\,m}=\widehat b^{\,n}\) for some \(m,n>0\), the selected
affine germs of \(A^m\) and \(B^n\) have the same base point and the same
positive slope, so they agree on a nondegenerate selected-side interval.
Both maps commute with the exact map \(R\); exactness transports this local
equality to the whole interval.  Thus \(A^m=B^n\), equivalently
\(P^{Lm}=Q^{Ln}\).  Lemma~\ref{lem:analytic-common-iterate} then gives a
common fixed point of \(P,Q\), contrary to the standing counterexample
assumption.  Therefore
\begin{equation}\label{term:eq:new-independent}
 \widehat a,\widehat b>1,\qquad
 \frac{\log\widehat a}{\log\widehat b}\notin\mathbb Q.
\end{equation}

\section{The unit-speed ray and a common Markov partition}

\subsection{Rebuilding the ray}

Using the positive right germ of the constant-slope map \(R\), define, for
\(t\geq0\),
\begin{equation}\label{term:eq:Psi}
 \Psi(t)=R^n\left(p_\theta+\frac{t}{\widehat s^{\,n}}\right)
\end{equation}
for any sufficiently large \(n\).  The definition stabilizes because
\[
 R(p_\theta+x)=p_\theta+\widehat s x
\]
on the selected germ.  One value of \(n\) works on every compact parameter
interval.  Thus \(\Psi\) is continuous and locally finite-piece PL.  On each
nonconstant affine piece its slope has absolute value
\(\widehat s^{\,n}/\widehat s^{\,n}=1\).

The ray is onto.  Indeed, choose a small selected-side interval contained in
the germ of \(R\).  Exactness sends its relative interior onto the whole
interval under some iterate, and \eqref{term:eq:Psi} places that image in the
range of \(\Psi\).  Near zero,
\begin{equation}\label{term:eq:Psi-initial}
 \Psi(t)=p_\theta+t.
\end{equation}
Commutation and the selected germs give
\begin{equation}\label{term:eq:Psi-equivariance}
\begin{split}
 \Psi(\widehat a t)&=A(\Psi(t)),\\
 \Psi(\widehat b t)&=B(\Psi(t)),\\
 \Psi(\widehat s t)&=R(\Psi(t)).
\end{split}
\end{equation}

\subsection{Lap-vertex values}

For an onto flat-free finite-piece PL map $T$, let
\[
 E(T)=\{0,1\}\cup\operatorname{Turn}(T),\qquad
 \mathcal V(T)=T(E(T)),
\]
where $\operatorname{Turn}(T)$ is the finite set of genuine changes of
monotonicity.

\begin{lemma}[Vertex values under composition]
\label{lem:vertex-composition}
For onto flat-free finite-piece PL maps $U,T$,
\begin{equation}\label{term:eq:vertex-composition}
 \mathcal V(U\circ T)=\mathcal V(U)\cup U(\mathcal V(T)).
\end{equation}
\end{lemma}

\begin{proof}
Let $x$ be a turn of the inner map $T$.  The two adjacent branches of $T$
approach $T(x)$ from the same side.  On that side $U$ has a nonzero affine
germ, even when $T(x)$ is itself a turn coordinate of $U$.  Thus the inner
turn cannot cancel in $U\circ T$.  Domain endpoints remain domain endpoints,
so
$U(\mathcal V(T))\subseteq\mathcal V(U\circ T)$.

Let $y$ be an interior turn coordinate of $U$.  Because $T$ is continuous
and onto, some preimage of $y$ is transverse: otherwise every preimage would
be a local extremum at level $y$, and no path in the domain joining a
preimage of $0$ to a preimage of $1$ could cross from one side of $y$ to the
other.  At a transverse preimage, $T$ is locally one-to-one and the outer
turn of $U$ produces a turn of the composite.  If $y$ is a range endpoint,
any point where $T$ assumes that global extremal value is either a domain
endpoint or a turn of $T$.  This accounts for $U(E(U))$ and proves
$\mathcal V(U)\subseteq\mathcal V(U\circ T)$.

Conversely, if $x\notin E(T)$, then $T$ is locally strictly monotone and
one-to-one at $x$; a turn of $U\circ T$ at $x$ can therefore occur only when
$T(x)$ is a turn of $U$.  If $x\in E(T)$, its composite value lies in
$U(\mathcal V(T))$.  These alternatives prove the reverse inclusion and
hence \eqref{term:eq:vertex-composition}.
\end{proof}

Define
\begin{equation}\label{term:eq:TA-TB}
 \mathcal T_A(m)=\bigcup_{0\leq i<m}A^i(\mathcal V(A)),\qquad
 \mathcal T_B(n)=\bigcup_{0\leq j<n}B^j(\mathcal V(B)).
\end{equation}
Iterating Lemma~\ref{lem:vertex-composition} and comparing $A^mB^n$ with
$B^nA^m$ gives, for all $m,n\geq1$,
\begin{equation}\label{term:eq:rectangle-identity}
 \mathcal T_A(m)\cup A^m(\mathcal T_B(n))
 =
 \mathcal T_B(n)\cup B^n(\mathcal T_A(m)).
\end{equation}

\subsection{Finiteness of the postvertex sets}

\begin{lemma}[Finite postvertex sets]
\label{lem:finite-postvertex}
The sets
\[
 \mathcal T_A=\bigcup_{i\geq0}A^i(\mathcal V(A)),\qquad
 \mathcal T_B=\bigcup_{j\geq0}B^j(\mathcal V(B))
\]
are finite.
\end{lemma}

\begin{proof}
Put
$\mathcal T=\bigcup_{m\geq1}\mathcal T_A(m)$ and
$\mathcal U=\mathcal V(B)$.  Taking $n=1$ in
\eqref{term:eq:rectangle-identity} gives
\[
 \mathcal T_A(m)\cup A^m(\mathcal U)
 =\mathcal U\cup B(\mathcal T_A(m)),
\]
and hence
\begin{equation}\label{term:eq:postvertex-differences}
\begin{aligned}
 \mathcal T_A(m)\setminus B(\mathcal T_A(m))&\subseteq\mathcal U,\\
 B(\mathcal T_A(m))\setminus\mathcal T_A(m)&\subseteq A^m(\mathcal U).
\end{aligned}
\end{equation}
If $y\in B(\mathcal T)\setminus\mathcal T$, choose $m_0$ so that a
$B$-preimage of $y$ belongs to $\mathcal T_A(m_0)$.  For every
$m\geq m_0$, the second inclusion in
\eqref{term:eq:postvertex-differences} yields $y\in A^m(\mathcal U)$.
Since $\mathcal U$ is finite, there are $u\in\mathcal U$ and $i<j$ with
$A^i(u)=A^j(u)=y$; therefore $A^{j-i}(y)=y$.  Such points lie on the finitely
many eventual cycles of the finitely many forward $A$-orbits starting in
$\mathcal U$.  Consequently
\begin{equation}\label{term:eq:B-T}
 B(\mathcal T)\subseteq\mathcal T\cup Z,
\end{equation}
where $Z$ is finite and every point of $Z$ is $A$-periodic.

Suppose $\mathcal T$ were infinite.  It is generated by the finite seed
$\mathcal V(A)$, so its points with infinite forward $A$-orbit fall into
finitely many tail-equivalence classes, where two points are equivalent when
their forward orbits meet.  If $w$ has infinite $A$-orbit, then $B(w)$ also
has infinite $A$-orbit.  Otherwise the identities
$A^iB(w)=B(A^i(w))$ would send infinitely many distinct points through the
finite-to-one map $B$ into a finite set.  Since the points of $Z$ have finite
orbits, \eqref{term:eq:B-T} makes $B$ a selfmap of the finite set of
infinite tail classes.

Choose a cyclic class of period $r\geq1$ and a representative $w$ in it.
Tail equivalence gives $p_0,q_0\geq0$ such that
\begin{equation}\label{term:eq:tail-relation}
 A^{p_0}B^r(w)=A^{q_0}(w).
\end{equation}
Commutation propagates this equality along the infinitely many distinct
points $A^i(w)$.  The two finite-piece PL maps
$A^{p_0}B^r$ and $A^{q_0}$ therefore agree at infinitely many points.  A
finite common affine partition has only finitely many boundary points, so
one open cell contains two, indeed infinitely many, of these agreement
points; the two affine restrictions agree on that cell.

Both maps commute with the exact map $R$.  If a nondegenerate agreement
interval $J_1$ satisfies $R^N(J_1)=\I$, then for $y=R^N(x)$ with
$x\in J_1$,
\[
 A^{p_0}B^r(y)=R^N A^{p_0}B^r(x)
              =R^N A^{q_0}(x)=A^{q_0}(y).
\]
Thus
\begin{equation}\label{term:eq:global-tail-relation}
 A^{p_0}B^r=A^{q_0}
\end{equation}
globally.  At every regular point, comparison of the constant absolute
slopes gives
$\widehat a^{p_0}\widehat b^r=\widehat a^{q_0}$, contrary to
\eqref{term:eq:new-independent}.  Hence $\mathcal T_A$ is finite.  The
same argument with $A,B$ interchanged proves that $\mathcal T_B$ is finite.
\end{proof}

\subsection{A common finite Markov set}

\begin{lemma}[Common finite Markov set]
\label{lem:common-Markov-set}
There is a finite set $\mathcal K\subseteq\I$ containing $0,1,p_\theta$
such that
\[
 A(\mathcal K)\subseteq\mathcal K,\qquad
 B(\mathcal K)\subseteq\mathcal K,
\]
and both $A$ and $B$ are affine and monotone between consecutive points of
$\mathcal K$.
\end{lemma}

\begin{proof}
Once $\mathcal T_A(m)$ has stabilized to $\mathcal T_A$,
\eqref{term:eq:postvertex-differences} applies for every larger $m$.  For
$u\in\mathcal V(B)$ whose forward $A$-orbit is eventually periodic, let
$C_A(u)$ denote its eventual cycle, and set
\[
 Z_A=\bigcup_{\substack{u\in\mathcal V(B)\\
                 u\text{ has eventually periodic }A\text{-orbit}}}C_A(u).
\]
The argument proving \eqref{term:eq:B-T} shows that
$B(\mathcal T_A)\setminus\mathcal T_A\subseteq Z_A$.  Thus $Z_A$ is finite
and consists of $A$-periodic points.  Choose $r_A$ divisible by the periods
of all points in $Z_A$.  The set
$\Fix(A^{r_A})$ is finite: on every affine cell the slope of $A^{r_A}$ has
absolute value $\widehat a^{r_A}>1$, so $A^{r_A}-\id$ cannot vanish on an
interval.  Both $A$ and $B$ preserve this fixed set.  Therefore
\[
 F_A=\mathcal T_A\cup\Fix(A^{r_A})
\]
is finite and forward invariant under both maps.  Construct $F_B$
symmetrically, using the eventual $B$-cycles met by the forward $B$-orbits
of $\mathcal V(A)$.

We next record turn transport.  Suppose $x$ is a turn of $A$ but not of $B$.
Then $B$ is locally affine and one-to-one at $x$.  In $BA$, the inner turn of
$A$ cannot cancel, because both branches approach $A(x)$ from the same side
and the corresponding germ of $B$ has nonzero slope.  Since $AB=BA$, the map
$AB$ also turns at $x$; local one-to-one-ness of $B$ then forces $B(x)$ to be
a turn of $A$.  If $x$ is a turn of both maps, then
$B(x)\in\mathcal V(B)\subseteq F_B$; the same conclusion holds for domain
endpoints.  Hence
\begin{equation}\label{term:eq:turn-transport}
 B(E(A))\subseteq E(A)\cup F_B,\qquad
 A(E(B))\subseteq E(B)\cup F_A.
\end{equation}
Applying the identical argument to $AR=RA$ and $BR=RB$ gives
\begin{equation}\label{term:eq:R-turn-transport}
 A(E(R))\subseteq E(R)\cup F_A,\qquad
 B(E(R))\subseteq E(R)\cup F_B.
\end{equation}

Set
\begin{equation}\label{term:eq:Markov-K}
 \mathcal K=
 F_A\cup F_B\cup E(A)\cup E(B)\cup E(R)\cup\{p_\theta\}.
\end{equation}
The sets $F_A,F_B$ are invariant under both maps;
$A(E(A))=\mathcal V(A)\subseteq F_A$ and
$B(E(B))=\mathcal V(B)\subseteq F_B$; and
\eqref{term:eq:turn-transport}--\eqref{term:eq:R-turn-transport} give all
remaining invariance checks.  The point $p_\theta$ is fixed by both maps.
Thus $\mathcal K$ is finite and forward invariant under $A,B$.  It contains
all turns of both maps.  Moreover, constant nonzero absolute slope excludes
an unrecorded same-direction kink: on two adjacent pieces with the same
monotonicity the signed slopes coincide, and continuity then makes the two
affine formulas identical.  Hence each map is affine and monotone on every
component of $\I\setminus\mathcal K$, and the endpoint images lie in
$\mathcal K$.
Since \(R=AB\), we also have \(R(\mathcal K)\subseteq\mathcal K\).
This is the required common Markov vertex set.
\end{proof}

\section{Common substitutions and Cobham periodicity}

\subsection{Directed edge substitutions}

Use both orientations of every \(\mathcal K\)-edge as a finite alphabet.
Since \(A\) is affine with nonzero slope on each edge and maps its endpoints
into \(\mathcal K\), the image of a directed edge is a nonempty directed edge
path.  This defines a nonerasing substitution \(\sigma_A\); define
\(\sigma_B\) similarly.  If a path description contains immediate
backtracking, retain the two oppositely directed letters rather than
cancelling them.

Because \(p<c<1\), we have \(p_\theta<1\).  Let \(k_+\) be the least member
of \(\mathcal K\) greater than \(p_\theta\), and put
\[
 e_+=(p_\theta,k_+),\qquad \ell=k_+-p_\theta.
\]
Because \(\mathcal K\) contains every turn of \(A,B,R\), and these maps have
constant absolute slope, their positive selected affine germs extend as
long as the argument remains in \(e_+\).  For \(0\leq t\leq\ell\), choose
\(N\) so large that \(t/\widehat s^{\,N}\) lies in the initial germ.
Every intermediate point
\(p_\theta+t/\widehat s^{\,j}\), \(0\leq j\leq N\), lies in the closure of
\(e_+\); iterating the positive branch of \(R\) in
\eqref{term:eq:Psi} therefore gives
\[
 \Psi(t)=p_\theta+t\qquad(0\leq t\leq\ell).
\]
Thus \(\Psi|_{[0,\ell]}\) parametrizes the directed edge \(e_+\) at unit
speed.  For every \(n\geq0\), applying
\eqref{term:eq:Psi-equivariance} \(n\) times gives
\[
 \Psi(\widehat a^{\,n}t)=A^n(\Psi(t)),\qquad
 \Psi(\widehat b^{\,n}t)=B^n(\Psi(t))
 \quad(0\leq t\leq\ell).
\]
Because \(\mathcal K\) is forward invariant under \(A,B\), the right-hand
sides traverse exactly the edge paths defining the corresponding
substitutions.  Hence, as literal parametrized paths,
\begin{equation}\label{term:eq:substitution-prefixes}
\begin{split}
 \sigma_A^n(e_+)
 &=\text{the itinerary of }\Psi\text{ on }[0,\widehat a^{\,n}\ell],\\
 \sigma_B^n(e_+)
 &=\text{the itinerary of }\Psi\text{ on }[0,\widehat b^{\,n}\ell].
\end{split}
\end{equation}
Since \(\widehat a,\widehat b>1\), both substituted words begin with
\(e_+\).  Their limits are the same one-sided directed-edge word \(x\), the
itinerary of the whole ray:
\begin{equation}\label{term:eq:common-fixed-word}
 \sigma_A(x)=x=\sigma_B(x).
\end{equation}

Discard every directed letter which never occurs in \(x\).  This leaves one
common finite alphabet invariant under both substitutions.  Indeed, if a
letter occurs in \(x\), every letter in its \(\sigma_A\)-image occurs in
\(\sigma_A(x)=x\), and the same is true for \(\sigma_B\).  The restrictions
are still nonerasing and prolongable on \(e_+\), and every retained letter
occurs in the displayed fixed point.

\subsection{Perron growth}

Use the standard incidence convention
$(M_A)_{e,f}=|\sigma_A(f)|_e$, and similarly for $B$: columns are source
letters and rows count output letters.  Assign to each directed geometric
edge its positive Euclidean length, and let $\boldsymbol\ell$ be the
resulting strictly positive column vector.  Since the image path of an edge
has Euclidean length equal to the absolute slope times the length of that
edge,
\begin{equation}\label{term:eq:length-eigenvectors}
 M_A^{\mathsf T}\boldsymbol\ell=\widehat a\boldsymbol\ell,\qquad
 M_B^{\mathsf T}\boldsymbol\ell=\widehat b\boldsymbol\ell.
\end{equation}

Recall that a \emph{Perron number} is a real algebraic integer
\(\lambda>1\) all of whose algebraic conjugates other than \(\lambda\) have
modulus strictly less than \(\lambda\).

\begin{lemma}[A weak Perron eigenvalue has a Perron power]
\label{term:lem:perron-power}
Let \(M\) be a finite nonnegative integral matrix.  If
\[
 Mv=\lambda v
\]
for a strictly positive vector \(v\) and a real number \(\lambda>1\), then
\(\rho(M)=\lambda\), and there is a positive integer \(u\) such that
\(\lambda^u\) is a Perron number.  Moreover \(M^u\) has spectral radius
\(\lambda^u\).
\end{lemma}

\begin{proof}
The Collatz--Wielandt inequalities applied to the strictly positive vector
\(v\) give \(\rho(M)=\lambda\).  Since \(M\) is integral, \(\lambda\) is an
algebraic integer; its minimal polynomial divides the characteristic
polynomial of \(M\), so every algebraic conjugate of \(\lambda\) is an
eigenvalue of \(M\) and has modulus at most \(\lambda\).

Put \(M\) in Frobenius block form.  If an eigenvalue \(z\) of \(M\) has
\(|z|=\lambda\), then it belongs to an irreducible diagonal block whose
spectral radius is \(\lambda\).  The periodic form of the Perron--Frobenius
theorem gives \(z=\lambda\zeta\), where \(\zeta\) is a root of unity.  Choose
\(u\) divisible by the orders of all the finitely many such \(\zeta\).

Let \(w\) be a conjugate of \(\lambda^u\).  An embedding of
\(\mathbb Q(\lambda^u)\) producing \(w\) extends to an embedding of
\(\mathbb Q(\lambda)\), so \(w=z^u\) for a conjugate \(z\) of \(\lambda\).
If \(|z|<\lambda\), then \(|w|<\lambda^u\).  If \(|z|=\lambda\), the
choice of \(u\) gives \(w=z^u=\lambda^u\).  Thus every distinct conjugate of
\(\lambda^u\) has modulus strictly less than \(\lambda^u\), which is exactly
the Perron property.  Finally
\(\rho(M^u)=\rho(M)^u=\lambda^u\).
\end{proof}

Apply Lemma~\ref{term:lem:perron-power} to
$M_A^{\mathsf T},M_B^{\mathsf T}$.  Transposition does not change the
spectrum.  There are positive integers $u,v$ such that
$\widehat a^{\,u},\widehat b^{\,v}$ are Perron numbers and are the spectral
radii of $M_A^u,M_B^v$, respectively.  The incidence matrix of a substitution
power is the corresponding matrix power.
Moreover
\[
 \sigma_A^u(x)=x=\sigma_B^v(x),
\]
the occurring alphabet is unchanged, and the substitutions remain
nonerasing and prolongable on \(e_+\).  Their iterated prefixes form
subsequences of the prefixes in \eqref{term:eq:substitution-prefixes}.  The
positive minimum and maximum edge lengths show that
\(\lvert\sigma_A^n(e_+)\rvert\) is comparable to \(\widehat a^{\,n}\), and
similarly for \(B\), so both prefix lengths tend to infinity.  Finally
\(\widehat a^{\,u}\) and \(\widehat b^{\,v}\) remain multiplicatively
independent, since a relation between positive powers of them would be a
relation between positive powers of \(\widehat a,\widehat b\).

\subsection{Durand's theorem and periodicity}

We use Durand's substitution form of Cobham's theorem
\cite[Definition 2 and Theorem 1]{Durand}.

\medskip
\noindent\textbf{Durand--Cobham theorem.}
\emph{Let \(\alpha,\beta>1\) be multiplicatively independent Perron
numbers.  Suppose a one-sided sequence over a finite alphabet is both
\(\alpha\)-substitutive and \(\beta\)-substitutive: more precisely, it is a
letter-to-letter coding of a fixed point of a substitution prolongable on a
generating letter whose occurring-letter incidence matrix has dominant
eigenvalue \(\alpha\), and it has an analogous presentation with dominant
eigenvalue \(\beta\).  Then the sequence is ultimately periodic.}
\medskip

Apply the theorem after replacing the substitutions by the powers obtained
above.  Each powered substitution is still prolongable on $e_+$ and fixes
$x$; every retained letter occurs in $x$ by construction; its incidence
matrix has dominant eigenvalue $\widehat a^{u}$ or $\widehat b^{v}$; and
the identity map is a letter-to-letter coding.  These are precisely the
conditions in Durand's Definition 2.  The two dominant eigenvalues are
multiplicatively independent Perron numbers, so the directed edge word $x$
is ultimately periodic.

Let \(\mathcal W\) be a repeated tail block of \(x\), and put
\begin{equation}\label{term:eq:period-length}
 P_0=\sum_{\text{occurrences }e\text{ in }\mathcal W}
          \operatorname{length}(e)>0.
\end{equation}
Every occurrence of a directed letter is traversed by \(\Psi\) at unit speed
from the same initial endpoint and in the same orientation.  Therefore
\begin{equation}\label{term:eq:eventual-metric-period}
 \Psi(t+P_0)=\Psi(t)
\end{equation}
for all sufficiently large \(t\), including block boundaries by continuity.

Equation \eqref{term:eq:eventual-metric-period} defines a continuous periodic
extension of the late tail to all of \(\mathbb R\): add a sufficiently large
integral multiple of \(P_0\) and evaluate \(\Psi\).  Eventual periodicity
makes the value independent of the choice.  Denote this extension by
\(\psi:\mathbb R\to\I\).  Its period group
\[
 \operatorname{Per}(\psi)
 =\{r\in\mathbb R:\psi(t+r)=\psi(t)\ \text{for every }t\in\mathbb R\}
\]
is a closed additive subgroup of \(\mathbb R\) containing the displayed
period \(P_0\).  It is not dense, since closedness would then make every real
number a period and \(\psi\) would be constant.  A nondense nonzero closed
subgroup of \(\mathbb R\) is generated by its least positive element.
Replace \(P_0\) by that element.  Thus the full period group is
\(P_0\mathbb Z\).
Equation \eqref{term:eq:Psi-equivariance} shows that
\(\widehat aP_0\) and \(\widehat bP_0\) are also periods.  Indeed, for all
sufficiently large \(t\), putting \(u=t/\widehat a\) gives
\[
 \psi(t+\widehat aP_0)
 =\psi\bigl(\widehat a(u+P_0)\bigr)
 =A(\psi(u+P_0))
 =A(\psi(u))
 =\psi(t).
\]
Periodic continuation makes this identity global; the argument for
\(\widehat b\) is identical.  Hence
\begin{equation}\label{term:eq:integer-slopes}
 \widehat a,\widehat b\in\mathbb N.
\end{equation}

Let \(J_{\rm tail}\) be the physical image of the periodic tail.  It is
nondegenerate and forward invariant under \(R\): for all sufficiently late
\(t\),
\[
 R(\Psi(t))=\Psi(\widehat s t)
\]
lies in the same tail image.  If \(J_{\rm tail}\) were proper, choose a
nonempty interval \(U\) contained in the relative interior of
\(J_{\rm tail}\) and open relative to the ambient interval \(\I\).  Every
\(R^n(U)\) would remain in \(J_{\rm tail}\), contradicting exactness of
\(R\).  Thus \(J_{\rm tail}=\I\).

The tail therefore defines an onto finite-piece PL unit-speed circle path
\begin{equation}\label{term:eq:psi-circle}
 \psi\colon\mathbb R/P_0\mathbb Z\longrightarrow\I.
\end{equation}
Using sufficiently late representatives, equation
\eqref{term:eq:Psi-equivariance} descends to
\begin{equation}\label{term:eq:psi-returns}
 \psi(\widehat a t)=A(\psi(t)),\qquad
 \psi(\widehat b t)=B(\psi(t)).
\end{equation}

\section{Original-phase affine lifts and the common fixed point}

\subsection{Cross-period globalization}

\begin{lemma}[Cross-period globalization]\label{term:lem:cross-period}
Let \(u,v\colon\mathbb R\to\I\) be nonconstant periodic finite-piece PL
surjections, all of whose nonzero slopes have absolute value one.  Suppose
that for some \(r>1\) and some interval map \(T\),
\[
 u(rt)=T(u(t)),\qquad v(rt)=T(v(t)).
\]
Then \(u\) and \(v\), up to a translation and possibly a reflection of the
parameter, are the same path.  In particular, they have the same full period
group and the same least positive period.
\end{lemma}

\begin{proof}
Choose a value $y$ which is not the image of a vertex of either polygonal
path.  Surjectivity supplies regular preimages $t_v,t_u$ of $y$.  Since the
nonzero slopes are $\pm1$, after shrinking a neighborhood there are
$\varepsilon\in\{1,-1\}$ and $\delta\in\mathbb R$ such that
\begin{equation}\label{term:eq:local-path-relation}
 v(t)=u(\varepsilon t+\delta)
\end{equation}
on a nonempty open interval $J$.

Iterating the two dilation identities shows that, on $J_n=r^nJ$,
\begin{equation}\label{term:eq:expanded-path-relation}
 v(t)=u(\varepsilon t+r^n\delta).
\end{equation}
Indeed, write $t=r^ns$ with $s\in J$ and apply $T^n$ to
\eqref{term:eq:local-path-relation}.  The lengths of the intervals $J_n$
tend to infinity.

Let $P_v,P_u$ be positive periods of $v,u$, and let $m_n$ be the midpoint of
$J_n$.  Reduce
\[
 \theta_n=m_n\pmod {P_v},\qquad
 \eta_n=\varepsilon m_n+r^n\delta\pmod {P_u}
\]
to fixed compact fundamental intervals.  Pass to a subsequence on which
$\theta_n\to\theta$ and $\eta_n\to\eta$.  For any fixed
$h\in\mathbb R$, the point $m_n+h$ belongs to $J_n$ for all sufficiently
large $n$.  Periodicity and \eqref{term:eq:expanded-path-relation} give
\[
 v(\theta_n+h)=u(\eta_n+\varepsilon h).
\]
Taking the limit yields
\[
 v(\theta+h)=u(\eta+\varepsilon h)
 \qquad(h\in\mathbb R).
\]
Thus, with $\delta_0=\eta-\varepsilon\theta$,
\[
 v(t)=u(\varepsilon t+\delta_0)\qquad(t\in\mathbb R).
\]
A number is a period of one side exactly when its signed counterpart is a
period of the other, so the two full period groups, and hence their least
positive periods, coincide.
\end{proof}

\subsection{Affine circle lifts of the original maps}

Recall that \(P,Q\) have constant absolute slopes \(C,D>1\), with
\begin{equation}\label{term:eq:root-slopes}
 C^L=\widehat a,\qquad D^L=\widehat b.
\end{equation}
Define
\[
 \psi_P(t)=P(\psi(t/C)).
\]
This is a unit-speed periodic surjection.  Since \(P\) commutes with \(A\),
equation \eqref{term:eq:psi-returns} gives
\[
 \psi_P(\widehat a t)=A(\psi_P(t)).
\]
Apply Lemma~\ref{term:lem:cross-period} to \(\psi_P,\psi\), with
\(r=\widehat a\) and \(T=A\).  They have the same least period \(P_0\).  But
\(CP_0\) is visibly a period of \(\psi_P\); hence \(C\) is a positive
integer.  Substituting \(t=Cx\) in the global path relation gives an affine
circle lift
\begin{equation}\label{term:eq:P-lift}
 P(\psi(x))=\psi(\mathcal P(x)),\qquad
 \mathcal P(x)=kx+\alpha\pmod{P_0},
\end{equation}
where \(k=\varepsilon_PC\in\mathbb Z\) and \(|k|\geq2\).  Similarly,
\begin{equation}\label{term:eq:Q-lift}
 Q(\psi(x))=\psi(\mathcal Q(x)),\qquad
 \mathcal Q(x)=\ell x+\beta\pmod{P_0},
\end{equation}
where \(\ell\in\mathbb Z\) and \(|\ell|\geq2\).

Physical commutation makes these affine lifts commute exactly.  Indeed,
\(\mathcal P\mathcal Q\) and \(\mathcal Q\mathcal P\) have the same linear
coefficient \(k\ell\), so their arguments differ by a constant \(\Delta\).
Their physical images under \(\psi\) agree.  Either affine composition is a
surjective circle covering, so
\[
 \psi(t+\Delta)=\psi(t)\qquad\text{for every }t.
\]
Minimality of \(P_0\) gives \(\Delta=0\pmod{P_0}\).

\subsection{The unique reflection}

\begin{lemma}[Globalization and uniqueness of the reflection]
\label{term:lem:global-reflection}
Let the same symbol \(\psi\) denote the \(P_0\)-periodic lift to
\(\mathbb R\) of the nonconstant unit-speed circle path
\eqref{term:eq:psi-circle}.  There is a unique reflection of the parameter
circle
\[
 \rho(t)=-t+\gamma\pmod{P_0}
\]
such that \(\psi\circ\rho=\psi\).
\end{lemma}

\begin{proof}
Choose an interior regular value \(y\): it lies in the interior of the range
of the path and is not the image of any vertex.  During one closed traversal,
the number of
upward crossings of \(y\) equals the number of downward crossings, and there
is at least one of each.  Choose real lifts \(t_+,t_-\) of an upward and a downward
crossing.  For some \(\varepsilon>0\),
\[
 \psi(t_++r)=y+r,\qquad
 \psi(t_-+r)=y-r
 \quad (|r|<\varepsilon).
\]
With
\[
 J_0=(t_+-\varepsilon,t_++\varepsilon),
 \qquad \gamma_0=t_++t_-,
\]
we therefore have
\begin{equation}\label{term:eq:local-reflection}
 \psi(t)=\psi(-t+\gamma_0)\qquad(t\in J_0).
\end{equation}

Iterating the first identity in \eqref{term:eq:psi-returns} gives, for every
\(n\geq0\),
\[
 \psi(\widehat a^{\,n}t)=A^n(\psi(t)).
\]
Apply this to both sides of \eqref{term:eq:local-reflection}.  On
\[
 J_n=\widehat a^{\,n}J_0,\qquad
 \gamma_n=\widehat a^{\,n}\gamma_0,
\]
it gives
\begin{equation}\label{term:eq:expanded-reflection}
 \psi(u)=\psi(-u+\gamma_n)\qquad(u\in J_n).
\end{equation}
Choose \(n\) so large that the length of \(J_n\) exceeds \(P_0\).  Given any
\(t\in\mathbb R\), choose \(m\in\mathbb Z\) with \(t+mP_0\in J_n\).
Periodicity and \eqref{term:eq:expanded-reflection} yield
\[
\begin{split}
 \psi(t)
 &=\psi(t+mP_0)
 =\psi(-(t+mP_0)+\gamma_n)\\
 &=\psi(-t+\gamma_n).
\end{split}
\]
Thus \(\rho(t)=-t+\gamma_n\) is a global reflection on
\(\mathbb R/P_0\mathbb Z\).  Choosing different real lifts of the two
crossings changes \(\gamma_0\) by an integral multiple of \(P_0\); since
\(\widehat a\in\mathbb N\) by \eqref{term:eq:integer-slopes}, it changes
\(\gamma_n\) by an integral multiple of \(P_0\), so the resulting circle
reflection is well-defined.

For uniqueness, suppose \(\rho_\gamma,\rho_\delta\) are two global reflection
symmetries.  Their composition is translation by \(\delta-\gamma\), so
\(\delta-\gamma\) is a period of \(\psi\).  The full period group of a
nonconstant continuous periodic function with least positive period \(P_0\)
is \(P_0\mathbb Z\): the period group is a closed subgroup of \(\mathbb R\),
and a nondiscrete one would be all of \(\mathbb R\), making the function
constant.  Hence \(\delta-\gamma\in P_0\mathbb Z\), and the two reflections
are identical on the parameter circle.
\end{proof}

Let
\begin{equation}\label{term:eq:reflection}
 \rho(t)=-t+\gamma,\qquad \psi\circ\rho=\psi
\end{equation}
be the unique reflection supplied by
Lemma~\ref{term:lem:global-reflection}.

We now prove that the lift \(\mathcal P\) preserves the \(\rho\)-orbits.
First,
\[
\begin{split}
 \psi(\mathcal P\rho(t))
 &=P(\psi(\rho(t)))\\
 &=P(\psi(t))
 =\psi(\mathcal P(t)).
\end{split}
\]
Write \(\mathcal P(t)=kt+\alpha\) and
\(\rho(t)=-t+\gamma\).  If \(y=\mathcal P(t)\), then, without choosing an
inverse branch,
\[
 \mathcal P\rho(t)=-y+k\gamma+2\alpha.
\]
Thus
\[
 \rho_P(y)=-y+k\gamma+2\alpha
\]
is another global reflection symmetry of \(\psi\): explicitly,
\(\mathcal P\rho(t)=\rho_P(\mathcal P(t))\), and every \(y\) equals
\(\mathcal P(t)\) for some \(t\), so \(\psi(\rho_P(y))=\psi(y)\).
Uniqueness gives \(\rho_P=\rho\), and therefore
\[
 \mathcal P\rho=\rho\mathcal P.
\]
The same argument gives \(\mathcal Q\rho=\rho\mathcal Q\).

Let
\[
 \pi\colon\mathbb R/P_0\mathbb Z\longrightarrow
 (\mathbb R/P_0\mathbb Z)/\langle\rho\rangle
\]
be the reflection quotient, which is an interval.  Since \(\psi\) is
constant on reflection orbits, there is a continuous surjection \(F\) with
\begin{equation}\label{term:eq:reflection-factor}
 \psi=F\circ\pi.
\end{equation}
Both affine lifts descend to commuting maps of this quotient interval.

\subsection{A common fixed reflection orbit}

Normalize the parameter circle to \(\mathbb R/\mathbb Z\) and translate so
that \(\rho(t)=-t\).  Write
\[
 \mathcal P(t)=kt+\alpha,\qquad
 \mathcal Q(t)=\ell t+\beta\pmod1.
\]
Commutation with \(\rho\) gives
\begin{equation}\label{term:eq:half-translations}
 \alpha,\beta\in\{0,\tfrac12\},
\end{equation}
while exact commutation of the two affine lifts gives
\begin{equation}\label{term:eq:affine-commutation}
 (k-1)\beta=(\ell-1)\alpha\pmod1.
\end{equation}
The following exhaustive parity check produces an orbit
\(\{t,-t\}\) fixed by both induced quotient maps.
\begin{enumerate}
 \item[(i)] If \((\alpha,\beta)=(0,0)\), take \(t=0\).
 \item[(ii)] If \((\alpha,\beta)=(\tfrac12,\tfrac12)\), then
 \eqref{term:eq:affine-commutation} says that \(k,\ell\) have the same
 parity.  Take \(t=\tfrac12\) if both are even and \(t=\tfrac14\) if both
 are odd.
 \item[(iii)] If \((\alpha,\beta)=(0,\tfrac12)\), then
 \eqref{term:eq:affine-commutation} says that \(k\) is odd.  Take
 \(t=\tfrac12\) if \(\ell\) is even and \(t=\tfrac14\) if \(\ell\) is odd.
 \item[(iv)] If \((\alpha,\beta)=(\tfrac12,0)\), then
 \eqref{term:eq:affine-commutation} says that \(\ell\) is odd.  Take
 \(t=\tfrac12\) if \(k\) is even and \(t=\tfrac14\) if \(k\) is odd.
\end{enumerate}
This calculation uses only parity and therefore includes negative degrees.
Let \(z=\pi(t)\) be the resulting quotient point.  It is fixed by both
induced maps.  Equations \eqref{term:eq:P-lift},
\eqref{term:eq:Q-lift}, and \eqref{term:eq:reflection-factor} therefore give,
for \(x=F(z)\),
\[
 P(x)=x=Q(x).
\]
This contradicts the standing assumption that \(P,Q\) have no common fixed
point and completes the proof of Theorem~\ref{thm:tame-C1}.
\hfill\(\square\)\par

\section*{Alignment with the formal proof}

The Lean development follows the proof above in the following dependency
order.  This list is also a guide to the exact interfaces between sections.

\begin{description}
\item[\texttt{Defs}, \texttt{DynamicsDefs}, \texttt{TameDefs}.]
The closed unit interval, iterates, fixed sets, one-sided germs, finite real
partitions, and the parameter-allowing o-minimal finite-piece \(C^1\)
predicate.

\item[\texttt{TameLemmas}, \texttt{QuotientDefs},
\texttt{QuotientLemmas}.]
Definability closure, one-variable o-minimal consequences, canonical finite
interval quotients, exact quotient fibres, respect by both generators, and
lifting of common fixed points.

\item[\texttt{Normalization}.]
The least-counterexample triple, finite fixed-component descent, selected
repelling directed state, permutation of the finite state set, and the
canonical return exponent.  The exponent is
\(2(2\#\Fix(W))!\), so it kills both the state permutation and its possible
orientation reversal.

\item[\texttt{FiniteArrival}.]
The labelled phase-intersection graph, interval Helly argument, dense
diagonal range, singleton endpoint fibre, strict endpoint-approach sequence,
interior-barrier basin closure, finite endpoint arrival, terminal-gap
exclusion, and local eventual onto.

\item[\texttt{OrbitDichotomy}, \texttt{CommonIterate}.]
The one-sided orbit line and orbit circle, local \(C^{1+BV}\) Kopell,
translation-number alternatives, Denjoy linearization, and the
Fugate--McLean reduction from equal positive powers to a common fixed point.

\item[\texttt{RayDynamics}.]
The canonical ray, tail core, dense-tail criterion, positive boundary lift,
logarithmic residue arithmetic, exactness, and the final disconnected-tail
quotient contradiction.

\item[\texttt{Polygonal}, \texttt{ConstantSlope},
\texttt{PolygonalRay}.]
The common Parry--Shultz coordinate, strict orientation of that coordinate,
constant-slope composition and iterate calculus, preservation of
multiplicative independence, explicit reconstruction of the positive
unit-speed polygonal ray, finite postvertex sets, and the common Markov set.

\item[\texttt{SymbolicDefs}, \texttt{SubstitutionPeriodicity}.]
Literal directed edges (with immediate backtracking retained), the two
substitutions fixing the same occurring-edge word, transpose-incidence length
eigenvectors, Perron powers, Durand--Cobham periodicity, metric eventual
periodicity, and the least-period return path.

\item[\texttt{OriginalPhase}, \texttt{AffineParity}.]
Cross-period globalization, expanding reflection globalization, normalized
integral affine lifts of the original maps, exact lift commutation, the
half-turn restriction on both translations, and the four parity cases.

\item[\texttt{Assembly}, \texttt{Main}.]
The two orbit alternatives are assembled into the contradiction and exposed
as the theorem that every commuting pair in the stated tame class has a
common fixed point.
\end{description}

Three interface corrections made visible by formalization are worth keeping
in the statement of any future version: the rebuilt ray needs a positive
right germ rather than only an absolute slope; the first Markov edge and the
postvertex argument need strict expansion of the relevant scales; and the
endpoint-basin lemma needs an interior fixed barrier rather than merely a
different fixed point.  Each strengthened hypothesis is proved before its
use in the main argument.

\begin{thebibliography}{99}

\bibitem{Boyce}
W.~M. Boyce,
\emph{Commuting functions with no common fixed point},
Trans. Amer. Math. Soc. \textbf{137} (1969), 77--92,
doi: \href{https://doi.org/10.1090/S0002-9947-1969-0236331-5}
{10.1090/S0002-9947-1969-0236331-5}.

\bibitem{Huneke}
J.~P. Huneke,
\emph{On common fixed points of commuting continuous functions on an
interval},
Trans. Amer. Math. Soc. \textbf{139} (1969), 371--381,
doi: \href{https://doi.org/10.1090/S0002-9947-1969-0237724-2}
{10.1090/S0002-9947-1969-0237724-2}.

\bibitem{AnusicMouron}
A.~Anu\v{s}i\'c and C.~Mouron,
\emph{Strongly commuting interval maps},
Fund. Math. \textbf{257} (2022), 39--68,
doi: \href{https://doi.org/10.4064/fm993-11-2021}
{10.4064/fm993-11-2021}.

\bibitem{Cohen}
H.~Cohen,
\emph{On fixed points of commuting functions},
Proc. Amer. Math. Soc. \textbf{15} (1964), 293--296,
doi: \href{https://doi.org/10.1090/S0002-9939-1964-0184219-1}
{10.1090/S0002-9939-1964-0184219-1}.

\bibitem{Folkman}
J.~H. Folkman,
\emph{On functions that commute with full functions},
Proc. Amer. Math. Soc. \textbf{17} (1966), 383--386,
doi: \href{https://doi.org/10.1090/S0002-9939-1966-0190916-6}
{10.1090/S0002-9939-1966-0190916-6}.

\bibitem{Joichi}
J.~T. Joichi,
\emph{On functions that commute with full functions and common fixed
points},
Nieuw Arch. Wisk. (3) \textbf{14} (1966), 247--251.

\bibitem{vanDenDries}
L.~van den Dries,
\emph{Tame Topology and O-Minimal Structures},
London Mathematical Society Lecture Note Series \textbf{248},
Cambridge University Press, Cambridge, 1998.

\bibitem{Navas}
A.~Navas,
\emph{Groups of Circle Diffeomorphisms},
Chicago Lectures in Mathematics,
University of Chicago Press, Chicago, 2011;
Propositions 2.2.2 and 2.2.3 and Theorems 3.1.1 and 4.1.1,
\url{https://arxiv.org/abs/math/0607481}.

\bibitem{Shultz}
F.~Shultz,
\emph{Dimension groups for interval maps II: the transitive case},
Ergodic Theory Dynam. Systems \textbf{27} (2007), 1287--1321;
Corollaries 4.4 and 4.6 and Proposition 3.6,
\url{https://arxiv.org/abs/math/0405467}.

\bibitem{Parry}
W.~Parry,
\emph{Symbolic dynamics and transformations of the unit interval},
Trans. Amer. Math. Soc. \textbf{122} (1966), 368--378,
doi: \href{https://doi.org/10.1090/S0002-9947-1966-0197683-5}
{10.1090/S0002-9947-1966-0197683-5}.

\bibitem{Durand}
F.~Durand,
\emph{Cobham's theorem for substitutions},
J. Eur. Math. Soc. \textbf{13} (2011), 1799--1814,
\url{https://ems.press/journals/jems/articles/4506}.

\bibitem{FugateMcLean}
J.~B. Fugate and T.~B. McLean,
\emph{Compact groups of homeomorphisms on tree-like continua},
Trans. Amer. Math. Soc. \textbf{267} (1981), 609--620,
doi: \href{https://doi.org/10.1090/S0002-9947-1981-0626493-5}
{10.1090/S0002-9947-1981-0626493-5}.

\end{thebibliography}

\end{document}
